\part{Module reference}\label{tm:part:module-reference}

\chapter{CPU core}\label{tm:ch:mod-cpu-core}

This chapter describes how each module in this area works. The area holds the pipeline, the register file and the control logic.

% --- generated by PR7c: begin
\section{\texttt{cpu/core/pipeline/rv32i\_pipeline\_ex1b.sv}}\label{tm:mod:rv32i-pipeline-ex1b}

This section holds comment text that was in the source file. The text is unchanged. Each line has no leading comment marker and no indentation.

\begin{relocated}{rtl/cpu/core/pipeline/rv32i_pipeline_ex1b.sv}

\subsection*{Header}

\begin{textblock}
Run-31: simplified field-pack stage — the 8-way case-mux has been removed.
EX1c now pre-decodes all trap outputs into dec_* fields registered in
ex1c_ex1b_reg_q.  EX1b reads those fields directly (1 gate each) instead
of re-computing from trap_type, eliminating the NOR5/NAND5 bottleneck
(~20 gate levels, 812 ps data arrival) that was the dominant timing violator
on the ex1c_ex1b_reg_q → EX2 path.
\end{textblock}

\end{relocated}

\section{\texttt{cpu/core/pipeline/rv32i\_pipeline\_ex1c.sv}}\label{tm:mod:rv32i-pipeline-ex1c}

This section holds comment text that was in the source file. The text is unchanged. Each line has no leading comment marker and no indentation.

\begin{relocated}{rtl/cpu/core/pipeline/rv32i_pipeline_ex1c.sv}

\subsection*{Header}

\begin{textblock}
Run-20: new purely-combinational stage inserted between ex1a_ex1b_reg_q
and ex1c_ex1b_reg_q to break the -782 ps timing critical path.

Run-35: restructured trap logic from 8-way if/else cascade to a
tree-balanced one-hot priority encoder + AND-OR mux.  No functional change.
The original design serialised: encoder (~7 gate levels) → decoder (~7 gate
levels) = ~14 gates deep.  The new design: parallel conditions (1-2 gates) →
balanced OR tree (3 gates) → one-hot sel (1 gate) → AND-OR mux (2 gates) =
~7 gates deep, targeting 720 ps / ~1389 MHz.
\end{textblock}

\subsection*{Before: ex1c\_o.dec\_trap\_valid      = 1'b0;}

\begin{textblock}
All data inputs (mtvec, irq_cause, pc, mret_pc, alu_result, immediate)
are registered: they are available at t=0 in this combinational stage.
The select (sel[]) resolves at ~4 gate levels; the AND-OR mux adds
~2 more levels.  Total: ~6 gate levels from registered inputs.

This path does NOT depend on trap_type — parallel to the encoder above.
\end{textblock}

\end{relocated}

\section{\texttt{cpu/core/pipeline/rv32i\_pipeline\_ex2.sv}}\label{tm:mod:rv32i-pipeline-ex2}

This section holds comment text that was in the source file. The text is unchanged. Each line has no leading comment marker and no indentation.

\begin{relocated}{rtl/cpu/core/pipeline/rv32i_pipeline_ex2.sv}

\subsection*{Header}

\begin{textblock}
This splits the former single EX combinational cone into two independent
timing windows, targeting ~200 ps WNS improvement on ASAP7.
\end{textblock}

\end{relocated}

\section{\texttt{cpu/core/pipeline/rv32i\_pipeline\_if.sv}}\label{tm:mod:rv32i-pipeline-if}

This section holds comment text that was in the source file. The text is unchanged. Each line has no leading comment marker and no indentation.

\begin{relocated}{rtl/cpu/core/pipeline/rv32i_pipeline_if.sv}

\subsection*{Before: assign dbg\_halted = dbg\_halt\_req \&\& !(ic\_valid\_o \&\& ic\_s...}

\begin{textblock}
Condition: !(ic_valid_o && ic_stall_i)
  When dbg_halt_req=1, ic_valid_o=0 (fetch is suppressed), so even if the
  I-cache has an in-flight AXI refill from before the halt arrived, the stall
  is for an abandoned request.  Blocking dbg_halted on ic_stall_i in that case
  deadlocks when the AXI handler has not yet started (e.g. ISA tests halt
  immediately after reset before starting AXI background tasks).
  Using ic_valid_o && ic_stall_i only stalls halt when the cache stall is for
  a live fetch request, preventing spurious halts mid-instruction.
\end{textblock}

\end{relocated}

\section{\texttt{cpu/core/rv32i\_alu.sv}}\label{tm:mod:rv32i-alu}

This section holds comment text that was in the source file. The text is unchanged. Each line has no leading comment marker and no indentation.

\begin{relocated}{rtl/cpu/core/rv32i_alu.sv}

\subsection*{Header}

\begin{textblock}
ADD/SUB/SLT/SLTU share one 33-bit adder, eliminating the redundant parallel
carry chains that were the binding setup-critical path (EX1a, Run 38 -3.89 ps).
\end{textblock}

\end{relocated}

\section{\texttt{cpu/core/rv32i\_control.sv}}\label{tm:mod:rv32i-control}

This section holds comment text that was in the source file. The text is unchanged. Each line has no leading comment marker and no indentation.

\begin{relocated}{rtl/cpu/core/rv32i_control.sv}

\subsection*{Before: always\_ff @(posedge clk or negedge rst\_n) begin}

\begin{textblock}
CRITICAL: We MUST register the branch/jump decision at the END of EXECUTE
state because:
1. In EXECUTE state, decoder outputs are valid for current instruction
2. When we transition to WRITEBACK (or MEM_WAIT), the instruction register
   latches the NEXT instruction, so decoder outputs change!
3. If we tried to use branch/jump signals combinationally in WRITEBACK,
   they would reflect the NEXT instruction, not the one we're committing!
\end{textblock}

\subsection*{Before: - Set when current\_state == HALTED and dbg\_step\_req is a...}

\begin{textblock}
Purpose: Carry the single-step intent through the pipeline.

The auto-clear logic in rv32i_cpu_top.sv clears dbg_ctrl_reg[2]
(and therefore dbg_step_req) on the HALTED→FETCH transition edge.
This means dbg_step_req is 0 by the time the CPU reaches WRITEBACK.
Without this flag the WRITEBACK→HALTED path would never be taken.

Protocol:
\end{textblock}

\end{relocated}

\section{\texttt{cpu/core/rv32i\_csr\_file.sv}}\label{tm:mod:rv32i-csr-file}

This section holds comment text that was in the source file. The text is unchanged. Each line has no leading comment marker and no indentation.

\begin{relocated}{rtl/cpu/core/rv32i_csr_file.sv}

\subsection*{Header}

\begin{textblock}
minstret / minstreth writes (bead f7z3): the write lands here in EX, but the
retire strobe trails EX by five stages, so instructions OLDER than the writer
are still in flight and the writer's own retirement is still to come.  Per the
privileged spec the older ones retire before the write and the writer's own
retirement does not increment, so none of them may land on the written value.
Each accepted write records, in a two-entry scoreboard, how many retire strobes
remain up to and including the writer (inflight_older_i + 1; the strobe of the
WB-stage instruction is the one in the write cycle itself and is overridden by
the write).  Those strobes are classified in program order:
  - the writer's own strobe: no increment at all (neither word);
  - older than a low-word (minstret) writer: no increment (overwritten);
  - older than a high-word (minstreth) writer: low word +1, but no carry into
    the high word it is overwriting.
Two entries cover the two writers that can be in flight together (EX-to-EX
spacing >= 3 cycles, retire 5 cycles after EX).  A trap entry does not touch the
scoreboard: the EX-stage trap point is two stages past EX, so no younger writer has
executed yet (spacing >= 3) and every pending writer is older than the trap and
still retires (an older writer's own strobe is at EX+5, never after the trap).
\end{textblock}

\subsection*{Before: pend    : entry holds a writer whose own retire strobe h...}

\begin{textblock}
=========================================================================
minstret write scoreboard (bead f7z3) -- see the header.
\end{textblock}

\subsection*{Before: logic       mrw\_pend0\_q, mrw\_pend1\_q;}

\begin{textblock}
drop / nocarry are registered so the retire -> counter enable sees flops only.
\end{textblock}

\subsection*{Before: 12'h7C0: begin  // dcache\_flush — CSRW dcache\_flush, x0 ...}

\begin{textblock}
csr_stall_i gate: when the EX stage is frozen (mem_cache_stall global
freeze), csr_access re-asserts every stall cycle because id_ex_reg is
held.  Normal CSR register writes are idempotent under re-fire, but
maintenance pulses are not — without this gate, the pulse would stay
high for the entire stall duration, causing the dcache to enter
CS_FLUSH_SCAN while the stalling store is still in CS_* → deadlock.
Fire exactly once: on the first (and only un-stalled) EX cycle.
\end{textblock}

\end{relocated}

\section{\texttt{cpu/core/rv32i\_hazard\_unit.sv}}\label{tm:mod:rv32i-hazard-unit}

This section holds comment text that was in the source file. The text is unchanged. Each line has no leading comment marker and no indentation.

\begin{relocated}{rtl/cpu/core/rv32i_hazard_unit.sv}

\subsection*{Before: logic load\_use\_ex1;}

\begin{textblock}
Run-20: the pipeline has THREE register stages between EX1a and MEM
(ex1a_ex1b_reg_q → EX1c comb → ex1c_ex1b_reg_q → EX1b comb →
 ex1_ex2_reg FF → MEM).  Three stall conditions are required:

  load_use_ex1   : load in EX1a (id_ex), consumer in ID (if_id)     → stall 1
  load_use_ex1c  : load in EX1c (ex1a_ex1b_reg_q, ex1b_ ports),
                   consumer still in ID                               → stall 2
  load_use_ex1b  : load in EX1b (ex1c_ex1b_reg_q, ex1c_ ports),
                   consumer still in ID                               → stall 3

After 3 stalls the consumer enters EX1a while the load is in MEM.
The D-cache stall (3 cycles) then holds the consumer in id_ex until
the load data appears in mem_wb_reg, which WB→EX1a forwarding delivers.
\end{textblock}

\end{relocated}

\section{\texttt{cpu/rv32i\_cpu\_top.sv}}\label{tm:mod:rv32i-cpu-top}

This section holds comment text that was in the source file. The text is unchanged. Each line has no leading comment marker and no indentation.

\begin{relocated}{rtl/cpu/rv32i_cpu_top.sv}

\subsection*{Before: always\_ff @(posedge clk\_i) begin}

\begin{textblock}
APB pready / prdata timing (GH #96 bead cyb):
  apb_prdata_o is a *registered* output (see always_ff below) — this
  breaks the ~800 ps combinational paddr[]/psel/pwrite -> prdata_o arc
  that the multi-clock SoC found inside this macro at the 780 ps cpu_clk
  period. Writes stay single-cycle (pready=1 immediately, matching prior
  behaviour exactly — dbg_ctrl_reg / breakpoint / GPR / PC writes are
  unaffected). Reads get exactly one APB3 wait state so the registered
  apb_prdata_o is valid on the same cycle apb_pready_o asserts.

  apb_pready_o is likewise a PURE registered output. It was previously
  `apb_pwrite_i ? 1'b1 : apb_read_wait_q`, and that one 2:1 mux exported a
  601.9 ps combinational arc from apb_pwrite_i in the macro Liberty —
  the same ~600 ps every APB input pays to cross this block's internal
  distribution, and the same reason apb_pslverr_o had to become a pure
  flop (see its comment below). Left combinational it would simply have
  become the next cpu_clk violator at SoC level after pslverr was fixed.
\end{textblock}

\subsection*{Before: Encoding (Doc: ds:mm:memory-map DBG\_STATUS bits [7:4]):}

\begin{textblock}
================================================================
Halt Cause Latch
================================================================
Capture the halt cause in a register so it is stable when DBG_STATUS
is read after the CPU enters HALTED.

Root problem: bp0_hit / bp1_hit are purely combinational and depend on
commit_valid_o, which is a single-cycle pulse from the FSM WRITEBACK
state.  By the time the CPU transitions to HALTED and the test reads
DBG_STATUS, commit_valid_o = 0, so bp0_hit = 0, and halt_cause = 0.

Fix: sample and hold the cause using commit_valid_o as the capture
trigger.  commit_valid_o is 1 exactly during WRITEBACK — the same cycle
bp0_hit / bp1_hit are valid.  We also cover the EBREAK case (already
handled by ebreak_halt) and the FETCH-time halt_req case.
\end{textblock}

\subsection*{Before: 1. EBREAK retires at WB (trap\_taken\_o \&\& trap\_cause\_o==3...}

\begin{textblock}
==================================================================
Priority (highest to lowest):
\end{textblock}

\end{relocated}
% --- generated by PR7c: end








\chapter{Memory and caches}\label{tm:ch:mod-memory-and-caches}

This chapter describes how each module in this area works. The area holds the instruction cache, the data cache and the memory blocks.

% --- generated by PR7c: begin
\section{\texttt{mem/rv32i\_cache\_pkg.sv}}\label{tm:mod:rv32i-cache-pkg}

This section holds comment text that was in the source file. The text is unchanged. Each line has no leading comment marker and no indentation.

\begin{relocated}{rtl/mem/rv32i_cache_pkg.sv}

\subsection*{Before: typedef enum logic [3:0] \{}

\begin{textblock}
Extended to 3 bits to accommodate synchronous-SRAM pipeline stages.
CS_HIT_PENDING added (Run 15 timing fix): one extra register stage between
CS_SRAM_LATCH and CS_TAG_CHECK so hit_comb is registered before it fans
out to the data SRAM mux network (~821 loads → registered hit_q/hit_bank_q).

Encoding rule (bead r51p, timing only, no behavioural change): state_q[3] = 1
exactly for the states that stall the pipeline in EVERY cycle they are occupied
(LATCH, HIT_PENDING, REFILL, WRITEBACK, UNCACHED_RD/WR, FLUSH_SCAN, INVAL_SCAN).
IDLE (stalls only with a request), TAG_CHECK (stalls only on a miss) and DONE
(never stalls) have state_q[3] = 0. dc_stall_o / ic_stall_o then need one state
bit instead of a 4-bit decode, which sat at the head of the Sky130 ss critical
cone (Doc: cr:defect:sky130-cpu-ss-setup-r51p). CS_IDLE stays 0 (reset state).
tb/tests/test_cache_stall_encoding.py guards the rule.
\end{textblock}

\end{relocated}

\section{\texttt{mem/rv32i\_clock\_gate.sv}}\label{tm:mod:rv32i-clock-gate}

This section holds comment text that was in the source file. The text is unchanged. Each line has no leading comment marker and no indentation.

\begin{relocated}{rtl/mem/rv32i_clock_gate.sv}

\subsection*{Header}

\begin{textblock}
Why: synthesis mapped the behavioural AND to an and2_2 (max_capacitance 0.19 pF at ss) and CTS then hung
~0.68 pF of root wire on it (the register-tree root buffer is placed at the sink centroid, ~4 mm from the
gate), failing CVT.MaxCapViolations at 8 corners. repair_design skips clock nets, so nothing else can fix
it; the driver has to be right from synthesis on. buf_16 has a 5.0 pF Liberty max_capacitance at ss (clkbuf_16: 1.017 pF
at ss, 1.53 pF at tt). It is a plain buffer, NOT a clkbuf_16, on purpose: clkbuf_16 is in CTS_CLK_BUFFERS, so TritonCTS
sees it as an existing clock buffer and skips the net ("CTS-0105: Net cpu_gated_clk already has clock buffer ...
Skipping...", RUN_2026-10-17_08-53-43: no clock tree built, setup WNS -130 ns). A gate whose output feeds no CTS tree (HARD_ROOT = 0) keeps the behavioural
\end{textblock}

\end{relocated}

\section{\texttt{mem/rv32i\_dcache.sv}}\label{tm:mod:rv32i-dcache}

This section holds comment text that was in the source file. The text is unchanged. Each line has no leading comment marker and no indentation.

\begin{relocated}{rtl/mem/rv32i_dcache.sv}

\subsection*{Before: CS\_HIT\_PENDING: combinational tag compare (tag\_dout\_r al...}

\begin{textblock}
Hit detection — pipelined (Run 15 timing fix)
\end{textblock}

\subsection*{Before: logic hit\_comb;}

\begin{textblock}
hit_bank_q[0:3]: per-bank registered copies for data SRAM write control.
Reduces post-synthesis fanout from ~821 to ~205 per copy.
\end{textblock}

\subsection*{Before: CS\_UNCACHED\_WR: begin}

\begin{textblock}
Completion: transition to CS_DONE (stall=0 for one cycle) rather
than directly to CS_IDLE.  Going straight to CS_IDLE would re-assert
dc_stall_o immediately (CS_IDLE && dc_valid_i=1) before the hazard
unit has a chance to advance the MEM stage, causing the same MMIO
address to be re-presented and re-entered.  CS_DONE holds stall=0
for one cycle so the pipeline can advance, then falls to CS_IDLE
with dc_valid_i already deasserted.
\end{textblock}

\end{relocated}

\section{\texttt{mem/rv32i\_icache.sv}}\label{tm:mod:rv32i-icache}

This section holds comment text that was in the source file. The text is unchanged. Each line has no leading comment marker and no indentation.

\begin{relocated}{rtl/mem/rv32i_icache.sv}

\subsection*{Before: logic addr\_mismatch;}

\begin{textblock}
After a branch redirect, pc_q changes while the cache may still be
processing the old (pre-redirect) address.  Detect this by comparing
the live fetch address (ic_addr_i = pc_q) against the latched address
(ic_addr_q) in the states that complete the previous SRAM lookup.
When a mismatch is detected the in-flight result is discarded and the
cache restarts for the new address.
\end{textblock}

\subsection*{Before: logic        tag\_write\_commit; // combinational commit q...}

\begin{textblock}
Tag SRAM write — registered commit (Run-17 P2 timing fix)

The combinational cone from state_q (CS_REFILL one-hot bit) into
tag_web0/tag_din0 had 298 fanout sinks and was the second-worst ASAP7
critical path at -606 ps.  Inserting a 1-cycle pipeline register
between the commit condition and the SRAM ports breaks this cone.

Latency impact: the tag SRAM write occurs 1 cycle after the last AXI
beat arrives instead of the same cycle.  The valid_array[idx] set is
delayed by the same 1 cycle (see valid_array always_ff below).
CS_REFILL transitions to CS_DONE (or CS_IDLE on mismatch) on the same
cycle as before — the registered write simply completes in the background
during CS_DONE, which is a no-stall cycle.  Correctness is preserved
because the CPU reads ic_rdata_o from refill_buf_q in CS_DONE, not from
the tag SRAM.
\end{textblock}

\subsection*{Before: logic idle\_read\_en;}

\begin{textblock}
SRAM read enable in CS_IDLE (bead r51p)

Functional arms: read only for a live fetch (ic_valid_i && !ic_invalidate_i).

Sky130 arm: read in every CS_IDLE cycle. A read issued while ic_valid_i = 0 (or during
FENCE.I) is never consumed: tag_dout_r / data_dout_r capture dout0 only in CS_SRAM_LATCH,
and CS_SRAM_LATCH is entered only from a CS_IDLE cycle with ic_valid_i = 1 and
ic_invalidate_i = 0, whose own read is the one that is captured. A read has no side
effect on the macro. In exchange csb0/addr0 no longer depend on ic_valid_i, i.e. on
the dcache-stall -> hazard-priority -> redirect/halt network that dominated the ss
setup failure of the CPU macro (Doc: cr:defect:sky130-cpu-ss-setup-r51p). Only that arm is
changed: on ASAP7 csb0 feeds the SRAM clock gate (extra enabled cycles = extra dynamic
power on a flow signed off as-is) and the other arms have no timing problem here.
Writes (tag_we_q / data_we_q) override below and are unaffected.
\end{textblock}

\subsection*{Before: logic        data\_write\_commit;}

\begin{textblock}
Data SRAM write — registered commit (Run-22 timing fix)

The combinational cone from state_q (CS_REFILL) + axi_rvalid_i into
data_csb0/web0/addr0/din0 had ~530+ fanout sinks and was the dominant
ASAP7 critical-path cluster after Run 21.  Inserting a 1-cycle pipeline
register between the commit condition and the SRAM ports breaks this cone,
mirroring the Run-17 P2 tag-SRAM registered-commit fix (tag_we_q above).

Latency impact: each refill beat's data SRAM write occurs 1 cycle after
the AXI beat arrives.  Beat 3 fires during CS_DONE (the no-stall cycle),
same timing as the tag write.  Correctness is preserved because the CPU
reads ic_rdata_o from refill_buf_q in CS_DONE, not from the data SRAMs.
valid_array is gated on tag_we_q, so the line is not visible as valid
until the tag write commits — by then all 4 data writes have also fired.
\end{textblock}

\end{relocated}
% --- generated by PR7c: end








\chapter{GPU}\label{tm:ch:mod-gpu}

This chapter describes how each module in this area works. The area holds the compute unit, the scheduler and the shared memory.

\chapter{SoC fabric}\label{tm:ch:mod-soc-fabric}

This chapter describes how each module in this area works. The area holds the crossbar, the bus bridges and the top-level wiring.

% --- generated by PR7a: begin
\section{\texttt{soc/apb\_cdc\_bridge.sv}}\label{tm:mod:apb-cdc-bridge}

This section holds comment text that was in the source file. The text is unchanged. Each line has no leading comment marker and no indentation.

\begin{relocated}{rtl/soc/apb_cdc_bridge.sv}

\subsection*{Header}

\begin{textblock}
Problem solved: soc_top.sv's APB_PLL2 slot (GH #92) sits between
axil_to_apb (an APB MASTER running on core_clk / the fabric domain) and
u_cpu_pll_sub's pll_apb_regs (an APB SLAVE running on cpu_clk_i / the
CPU-domain PLL reference — a genuinely different clock once GH #93
re-sources the CPU from cpu_core_clk). apb_interconnect in between is
purely combinational (no clk port, see soc_bus.sv "5. APB interconnect"),
so without this bridge the psel/penable/pwrite/paddr/pwdata/pstrb request
and the pready/prdata/pslverr response cross a clock boundary with zero
synchronisation — a live, unwaived CDC hazard the moment cpu_clk_i !=
clk_i (soc_top.sv's now-removed ⚠️ block documented this and deferred the
fix to GH #93/#95). GH #95 reused this same module for the APB3 debug
port (u_apb_dbg_cdc) too.

── Why 4-phase, not a documented contract (bead claude_verilog_test-rfz) ──
The ORIGINAL version of this module used a 2-phase (NRZ) toggle handshake:
a single bit flipped once per transfer, and the peer domain detected a
transfer by comparing its synchronised copy against a locally-tracked
"last seen" value (req_seen_q / ack_seen_q, now removed — see below). PULP
(cdc_reset_ctrlr.sv) and OpenTitan (prim_sync_reqack.sv) both
independently conclude that encoding is unsafe under ONE-SIDED
asynchronous reset, and both ship a 4-phase RZ variant specifically to
fix it:
  - PULP: "We use a 4-phase rather than a 2-phase CDC to avoid the issues
    of one-sided async reset that might trigger spurious transactions."
  - OpenTitan: its RZ option "is safe to reset either domain in isolation,
    since the two FSMs cannot get out of sync due to persistent EVEN/ODD
    states." Its NRZ default instead carries an explicit machine-checked
    contract that resetting one domain requires resetting the other AT
    THE SAME TIME.
Three independent async-APB bridges surveyed for this bead — ARM
cmsdk_ahb_to_apb_async.v, Hazard3 hazard3_apb_async_bridge, and OpenDAP —
are all 4-phase, and OpenDAP states the failure mode in exactly these
terms: "A NRZI toggle handshake ... can cause spurious bus accesses when
only one side of the link is reset." This module's structure already
matched ARM's bridge exactly (2x 2-FF sync on the handshake only, payload
held stable and read combinationally once qualified) — only the encoding
differed, and NRZ was the one all three surveyed projects avoided.

The cheaper alternative — document + assert a contract that a CPU-domain
PMU reset must not overlap an in-flight APB_PLL2 transaction — was
evaluated and REJECTED: OpenTitan's NRZ contract requires that resetting
one domain also resets the other, OVERLAPPING — which this SoC's PMU
cannot honour, since PD_CPU power-down resets the CPU domain (this
module's m_* face for u_apb_pll2_cdc) ALONE by design. Worse, the PMU
power-down path has no drain/quiesce step (a known, documented gap —
see soc_top.sv), so nothing guarantees the config path is idle before a
one-sided reset can land. Documenting a contract the design already
violates is not a fix.

GH #95 directed testing (test_reset_asym_s_held_m_free,
test_reset_asym_staggered_s_then_m, test_reset_asym_staggered_m_then_s —
3 stagger windows x 2 release orders each) found NO counterexample under
the OLD NRZ encoding. This conversion is therefore a structural fix
against a class two upstream projects and three independent bridges treat
as real, not a response to an observed simulation failure. Throughput is
halved (two round trips instead of one per transfer, see protocol below);
this is fine and expected on a config/debug path (APB_PLL2, APB3 debug),
never a data path.
\end{textblock}

\subsection*{Before: input  logic               scanmode\_i,}

\begin{textblock}
── DFT scan bypass (bead claude_verilog_test-07n) — global, not
   per-face: forwarded to BOTH internal cdc_reset_sync instances
   (u_s_rst_sync, u_m_rst_sync). Must be connected explicitly at
   every instantiation (no SV port defaults); tie scanmode_i=1'b0,
   scan_rst_ni=1'b1 where no scan flow exists — see
   cdc_reset_sync.sv for the full rationale. ─────────────────────────
\end{textblock}

\subsection*{Before: logic m\_in\_rst\_s;   // 1 = destination domain observed i...}

\begin{textblock}
── Destination-reset status, synchronised INTO the source domain ──────
Single-bit level crossing of the raw m_rst_n_i input (not m_rst_sync_n
— we want the source domain's own view of "is the destination
currently reset", independent of the destination's internal sync
latency) via the same cdc_sync primitive already used for the
req/ack levels.

Polarity is DELIBERATE: we synchronise the ACTIVE-HIGH "destination in
reset" sense (~m_rst_n_i), not m_rst_n_i directly, so that
cdc_sync's own fixed reset-to-0 behaviour means "presumed NOT in
reset" (i.e. presumed alive) the instant s_rst_sync_n releases. This
is not just a style choice: an earlier version of this synchroniser
sampled m_rst_n_i directly (reset-to-0 = "presumed IN reset"), which
is conservative for the availability fix in isolation but silently
STACKS a second synchroniser's settling latency onto the very first
transaction after every s_rst_n_i release — even when m_rst_n_i was
never actually asserted — because the reset-default (0) disagreed
with the immediate steady-state truth (1/alive), and it takes
SYNC_STAGES_TO_S more s_clk_i cycles for the real value to arrive. In
the s_ns=18/m_ns=4 clock-ratio regression this cost exactly one
spurious pslverr=1 on the first post-reset write. With the inverted
polarity, the reset default (0 = "not in reset") already MATCHES
steady-state truth whenever m_rst_n_i is genuinely high, so the good
path incurs no extra latency. The bad path (m_rst_n_i genuinely low)
is unaffected: the destination's own req_m/ack_q/d_state_q stay
async-held at 0 by m_rst_sync_n for as long as m_rst_n_i is actually
low (rst_both_n includes it), so a request forwarded during the few
cycles before this synchroniser catches up is invisible to the
destination and is cleaned up (req_q clamped to 0, S_REQ aborted with
pslverr) the moment m_in_rst_s does catch up — still bounded, still no
fabricated write, just detected a few cycles later in that narrow
race. See header "Reset strategy".
\end{textblock}

\subsection*{Before: logic m\_in\_rst\_s;   // 1 = destination domain observed i...}

\begin{textblock}
Why documentation, not enforcement: the reviewer's proposed
alternative (gate s_psel_i && !s_penable_i with an asserted reset
value) is exactly the conservative polarity already tried and
rejected — it regressed test_clock_ratio_slow_s_fast_m at an 18 ns/
4 ns clock ratio by stacking a full extra SYNC_STAGES_TO_S-cycle
settle latency onto the FIRST post-reset transaction even when
m_rst_n_i was never actually asserted (one spurious pslverr=1 that
does not exist today). The forwarding decision (req_fwd_q) is
already the narrowest thing this status bit gates — the SETUP is
still accepted into S_REQ unconditionally; only the forward is
gated — so there is no cheaper enforcement point to move the gate
to. The only way to also close the window itself (rather than just
bound and no-op it, as done here) is a synchroniser whose reset
default depends on knowing m_rst_n_i's true value at reset time,
which is precisely what a synchroniser resolves only after
SYNC_STAGES_TO_S clock edges — not achievable without deviating from
the shared cdc_sync primitive and its per-domain-synchronised-
reset discipline used uniformly everywhere else in this file. Given
the window is already proven bounded and side-effect-free by (a)-(c),
that deviation is not worth the extra risk. Documentation is correct.
\end{textblock}

\subsection*{Before: rst\_both\_n still exists and still couples the DESTINATIO...}

\begin{textblock}
── Reset: ASYMMETRIC — see "Reset strategy" in the module header ───────
\end{textblock}

\subsection*{Before: strictly required for RZ correctness but retained for th...}

\begin{textblock}
BOTH resets (a conservative, reviewed choice — see header, no longer
\end{textblock}

\subsection*{Before: logic              req\_fwd\_q;}

\begin{textblock}
header "Reset strategy" / "Availability fix".
\end{textblock}

\end{relocated}

\section{\texttt{soc/apb\_interconnect.sv}}\label{tm:mod:apb-interconnect}

This section holds comment text that was in the source file. The text is unchanged. Each line has no leading comment marker and no indentation.

\begin{relocated}{rtl/soc/apb_interconnect.sv}

\subsection*{Before: localparam int unsigned SEL\_W = \$clog2(N\_SLAVES + 1);}

\begin{textblock}
── Address decode ───────────────────────────────────────────────────────
sel_idx: index of the selected slave, or N_SLAVES when no match.
First-match (lowest index wins) — iterate up and stop on first hit.

Two widths are needed and they are NOT the same number:
  SEL_W — wide enough to also hold the N_SLAVES "no match" sentinel.
  IDXW  — wide enough to index the [N_SLAVES] port arrays, and no wider.
They differ exactly when N_SLAVES is a power of two (e.g. N_SLAVES=8:
SEL_W=4, IDXW=3), which is what made the response mux trip WIDTHTRUNC
when Phase 6a took the APB sub-tree from 7 slaves to 8. The `> 1` guard
on IDXW mirrors apb4_register_bank.sv's own IDXW localparam and keeps
N_SLAVES==1 (this module's default) legal — $clog2(1) is 0, and a bare
$clog2(N_SLAVES)-1:0 slice would be [-1:0].
\end{textblock}

\end{relocated}

\section{\texttt{soc/async\_axi\_fifo.sv}}\label{tm:mod:async-axi-fifo}

This section holds comment text that was in the source file. The text is unchanged. Each line has no leading comment marker and no indentation.

\begin{relocated}{rtl/soc/async_axi_fifo.sv}

\subsection*{Before: parameter int unsigned AW    = axi\_pkg::AXI\_ADDR\_WIDTH,}

\begin{textblock}
── Bus geometry (defaults from axi_pkg, scope-resolved — no file-scope
   import; see rtl/soc/axi4_crossbar.sv:40-46) ─────────────────────────
\end{textblock}

\subsection*{Before: input  logic               scanmode\_i,}

\begin{textblock}
── DFT scan bypass (bead claude_verilog_test-07n) — global, not
   per-face: forwarded to BOTH internal cdc_reset_sync instances
   (u_s_rst_sync, u_m_rst_sync). Must be connected explicitly at
   every instantiation (no SV port defaults); tie scanmode_i=1'b0,
   scan_rst_ni=1'b1 where no scan flow exists — see
   cdc_reset_sync.sv for the full rationale. ─────────────────────────
\end{textblock}

\end{relocated}

\section{\texttt{soc/axi4\_crossbar.sv}}\label{tm:mod:axi4-crossbar}

This section holds comment text that was in the source file. The text is unchanged. Each line has no leading comment marker and no indentation.

\begin{relocated}{rtl/soc/axi4_crossbar.sv}

\subsection*{Before: logic [N\_SLAVES-1:0][IW-1:0]    \_s\_awid;}

\begin{textblock}
─────────────────────────────────────────────────────────────────────────
Intermediate packed wires for all driven output port arrays.
always_comb drives these (Synlig-safe: internal wires, not ports).
Continuous assign in generate connects them to the unpacked port arrays.
Synlig/UHDM bug: always_comb driving port array slices via loop index
produces empty wire names in RTLIL (kernel/rtlil.cc:2150 assert).
Continuous assign inside generate is not affected by this bug.
─────────────────────────────────────────────────────────────────────────
Slave-side output intermediates (N_SLAVES slices)
\end{textblock}

\end{relocated}

\section{\texttt{soc/pmu.sv}}\label{tm:mod:pmu}

This section holds comment text that was in the source file. The text is unchanged. Each line has no leading comment marker and no indentation.

\begin{relocated}{rtl/soc/pmu.sv}

\subsection*{Header}

\begin{textblock}
Power-up needs NO ack for the iso-release-vs-reset-release order (decision, bead
3ogd). The UPF's last up-step, de-assert iso, stays last, and the hazard on the way
up is the mirror of the one on the way down: iso released too EARLY. It is harmless
here by structure, for any ratio: the only things the clamp holds are lines sampled
by cpu_gated_clk (u_cpu_axi_cdc's s-side, the CPU) plus the commit_valid_o
observation pin, and cpu_gated_clk cannot tick until the reset release is observed
in the CPU domain (soc_top: cpu_gated_clk_en includes pmu_cpu_rst_n_cpu_sync). A
clamp released a cpu_core_clk edge ahead of the reset release (a metastable first
stage in silicon; in this zero-delay model both paths are 2-flop syncs on one clock
and keep the nominal order) therefore exposes only retained values on a still-stopped
clock. Way down, by contrast, the reset acts ASYNCHRONOUSLY and at once, so the
clamp being late is observable: that direction gets the ack. (The release ack above
is a handshake-completion wait, not an ordering wait.)
\end{textblock}

\subsection*{Header}

\begin{textblock}
Reset style (explicit decision, see Doc: um:guide:coding-guidelines
section 1.4): this module uses SYNCHRONOUS reset (`always_ff @(posedge
clk_i)` gated by `if (!rst_n_i)`), matching every other module on the APB
peripheral ring it lives on (uart_controller, spi_controller, timer,
interrupt_controller, apb4_register_bank, pll_apb_regs — all synchronous-
reset, PD_PERIPH always-on domain). Section 1.4's general guidance for NEW
modules is async-assert/sync-deassert with a per-domain reset synchronizer,
but soc_top has NO reset synchronizer anywhere in the design (core_rst_n =
rst_n_i & pll_locked is a plain combinational AND — pll_subsystem.sv:93) and
section 1.4 also requires uniform discipline within one subsystem. Adding
this module's own async-reset flavor (and the synchronizer that would then
be needed to make it safe) would make the PMU the sole exception on its own
bus and does not match the surrounding synchronous-reset peripheral
subsystem it integrates into. Matching the subsystem was chosen over
introducing the project's first reset synchronizer as a side effect of a
PMU; full CPU/peripheral reset-discipline unification remains the separate
P3 audit item already tracked in Doc: um:guide:coding-guidelines.
\end{textblock}

\subsection*{Before: logic clk\_en\_q      [N\_DOM];}

\begin{textblock}
Registered sequencing outputs (ktwt). These drive cdc_sync D pins and the
ASYNC clear of u_cpu_pmu_rst_sync in soc_top, and (GPU) the async reset
and the handshake isolation mux of the GPU domain. A combinational
decode of the 4-bit binary dom_state_q can glitch when several state
bits change on one edge (4 of the 9 transitions; a runt reset on
DOM_PU_RESTORE -> DOM_PU_ISO), and a glitch on an async clear or a
synchroniser D pin is a real event. Each output is therefore the Q of
its own flop, loaded with the decode of the NEXT state: the value seen
at cycle t is exactly decode(dom_state_q at t), so no output moves one
cycle earlier or later than in the combinational version (UPF order and
test_pmu's per-cycle expectations are unchanged).
\end{textblock}

\subsection*{Before: parameter int unsigned ISO\_ACK\_TIMEOUT = 255}

\begin{textblock}
STATUS[3] (see "Isolation handshake"). >= 2; elaboration guard below.
\end{textblock}

\subsection*{Before: DOM\_PU\_ISO:     dom\_state\_d[d] = (!iso\_ack\_w[d] || wait\_...}

\begin{textblock}
Hold until the release is acknowledged (4-phase RZ; see header).
\end{textblock}

\subsection*{Before: Reset values are decode(DOM\_ON): ungated, not isolated, ...}

\begin{textblock}
── State + output registers (synchronous reset -- see header) ──
\end{textblock}

\end{relocated}

\section{\texttt{soc/soc\_top.sv}}\label{tm:mod:soc-top}

This section holds comment text that was in the source file. The text is unchanged. Each line has no leading comment marker and no indentation.

\begin{relocated}{rtl/soc/soc_top.sv}

\subsection*{Before: logic pmu\_cpu\_clk\_dis\_sync;}

\begin{textblock}
Revised ordering (cold-boot / PLL-(re)lock, the scenario this bug
report is about): while cpu_core_rst_n is low, pmu_cpu_clk_dis_sync is
held at its async-reset value of 0 (not disabled), so
cpu_gated_clk = cpu_core_clk unconditionally for the ENTIRE duration
cpu_core_rst_n (and therefore cpu_domain_rst_n, which can only release
no earlier than cpu_core_rst_n) is low — i.e. as many cpu_gated_clk
edges as cpu_core_rst_n's assertion lasts (STUB_LOCK_CYCLES=16 in the
PLL stub, or the real PLL's lock time), far more than the 1 edge
u_cpu's synchronous reset needs to load RESET_PC and its other reset
values. cpu_core_rst_n releases; pmu_cpu_rst_n_cpu_sync is already 1
(per the analysis above); cpu_domain_rst_n releases on that same edge,
and cpu_gated_clk has been ticking every cycle before, through, and
after that edge — no gap, so the release edge itself is also cleanly
observed by u_cpu. Only afterward, once firmware ever asserts
pmu_cpu_clk_en=0, does pmu_cpu_clk_dis_sync begin tracking it (2-cycle
pipeline latency), which is the normal, intended clock-gating path.
\end{textblock}

\subsection*{Before: logic pmu\_cpu\_clk\_dis\_sync;}

\begin{textblock}
GH #90 bead 542 fix (live PMU power-cycle ordering hazard, found during
the GH #93 sanity check above and deliberately deferred at the time):
pmu.sv's per-domain sequencer asserts rst_n_w=1 one full core_clk cycle
BEFORE it asserts clk_en_w=1 on the power-UP path (DOM_PU_RST releases
reset; DOM_PU_CLK, the next state, ungates the clock) — and mirrors
that on power-DOWN (clk_en_w drops one cycle before rst_n_w drops).
This ordering is NOT a pmu.sv implementation choice: docs/
Doc: tm:pd:power-domain-evaluation:3-what-is-doable-now-the-scoped-go and pmu.sv's own header fix this
order as the UPF-mandated sequence, and it is the mechanism that lets
this behavioral model preserve CPU architectural state (PC/MSTATUS/
MTVEC) across an off->on cycle WITHOUT real retention-cell silicon: as
long as u_cpu's clock never ticks while its domain reset is asserted,
its synchronous `if (!rst_n)` branch never executes during a PMU
power-cycle, so its flops simply hold whatever value they had before
the domain went quiet. Editing pmu.sv's state order to fix this bead
would either violate that documented UPF sequencing or defeat the
retention-by-construction trick (see pmu.sv header) — so this fix is
deliberately in soc_top.sv, not pmu.sv. (That trick is CPU-only and a
behavioural-model artifact: the GPU domain is async-reset, so a PMU
power-cycle DOES clear GPU state — bead 4a7i, pmu.sv header.)
\end{textblock}

\subsection*{Before: logic pmu\_cpu\_clk\_dis\_sync;}

\begin{textblock}
Worst-case skew tolerance: UNBOUNDED. The fix does not rely on
bounding, matching, or reasoning about the relative latency of the two
independent synchronisers at all — it makes cpu_domain-observed PMU
reset the master gating term for the CPU's own clock, so no amount of
relative skew between u_cpu_pmu_rst_sync and u_cpu_clk_dis_sync can
produce a clock edge while that reset is asserted. It is symmetric
across both directions: on power-up, the gate cannot open until
pmu_cpu_rst_n_cpu_sync is already 1; on power-down, the gate closes the
instant pmu_cpu_rst_n_cpu_sync drops to 0, even if
pmu_cpu_clk_dis_sync has not yet caught up to reflect clk_en_w=0.
\end{textblock}

\subsection*{Before: logic pmu\_cpu\_iso\_en\_sync;}

\begin{textblock}
Safe-reset-value sanity check (GH #93-fix review): unlike the clk_en
crossing above, this one is NOT inverted, and that is correct.
cdc_sync resets pmu_cpu_iso_en_sync to 0 = "not isolated", which
matches pmu_cpu_iso_en's own PMU-side reset default (DOM_ON: iso_en_w=0,
see pmu.sv reset-defaults comment) — the synchroniser's safe state and
the source signal's safe state agree, so there is no cold-boot window
where this crossing disagrees with its source the way the un-inverted
clk_en sync did. (Whether "not isolated" is itself the physically safe
choice for isolation cells in general is a separate, pre-existing
design decision — see the ISO_CPU clamp-scope note below — not one this
fix revisits.)
\end{textblock}

\subsection*{Before: logic pmu\_cpu\_rst\_n\_cpu\_sync;}

\begin{textblock}
Ordering sanity check (GH #93-fix review): rst_n_i here is deliberately
pmu_cpu_rst_n, NOT cpu_core_rst_n. cdc_reset_sync's rst_n_i is the
SOURCE async reset event being synchronised into clk_i's domain (see
cdc_reset_sync.sv header and its other two uses inside
async_axi_fifo.sv, which feed it s_rst_n_i/m_rst_n_i — the real per-
domain reset events, never a foreign domain's reset). cpu_core_rst_n is
already NATIVE to the cpu_core_clk destination domain (produced by
u_cpu_pll_sub, synchronous to cpu_core_clk already), so it needs no
synchroniser of its own; it is combined with pmu_cpu_rst_n_cpu_sync by
the plain AND below instead. Feeding cpu_core_rst_n into this
synchroniser's rst_n_i in place of pmu_cpu_rst_n would be wrong (it
would synchronise the wrong signal) and feeding it in ADDITION (e.g. as
an extra async clear) would be redundant — cpu_domain_rst_n's
combinational AND already asserts whenever EITHER source is low. No
ordering surprise: this is the correct, intentional use of the
primitive.
\end{textblock}

\subsection*{Before: logic uart\_irq;}

\begin{textblock}
Declared here (ahead of u_ext_irq_sync / u_timer_irq_sync below, which
read ext_irq/timer_irq as CDC synchroniser inputs) rather than further
down with the rest of the IRQ signal group, so every use is strictly
after its declaration (yosys-slang strict-mode requirement; no
functional change — see bead q7n).
\end{textblock}

\subsection*{Before: logic cpu\_domain\_rst\_n, gpu\_domain\_rst\_n;}

\begin{textblock}
GPU/fabric domain: left UNCHANGED (core_rst_n & pmu_gpu_rst_n, no new
reset synchroniser added). Rationale: the GPU/bus/peripherals all still
share the single core_clk domain (only the CPU moved), so there is no
new clock boundary here to synchronise across, and gpu_domain_rst_n's
existing timing has 142/142 soc_all regression history behind it.
Inserting a cdc_reset_sync here would add STAGES cycles of reset-
deassert latency to every peripheral for zero CDC benefit and is
deliberately NOT done, to avoid perturbing that regression baseline.
\end{textblock}

\subsection*{Before: logic        dbg\_bridge\_m\_psel, dbg\_bridge\_m\_penable, db...}

\begin{textblock}
Fix: a second apb_cdc_bridge instance, u_apb_dbg_cdc, restores the
pre-#93 external contract ("drive apb_* synchronously to clk_i") by
bridging INTO the CPU domain instead of redefining that contract.
Modelled on ARM CoreSight SoC-400's DAPBUS async bridge (DDI 0480F
§4.2.6/§4.2.8): the debug port clock is asynchronous to the block it
drives by construction (an external probe cannot guarantee edge
alignment with cpu_clk_i), and the async bridge is what lets debug stay
power-isolated from the domain it accesses.
\end{textblock}

\subsection*{Before: async\_axi\_fifo u\_cpu\_axi\_cdc (}

\begin{textblock}
RETAIN vs FLUSH — the actual bug this rewiring closes: a PMU-driven
live CPU-domain power-cycle must RETAIN u_cpu's architectural state
(that is the entire point of clock-gate-only power-down — see
pmu_cpu_rst_n_cpu_sync's ordering note above u_cpu_cg: u_cpu is
synchronous-reset and cpu_gated_clk is forbidden from ticking while
that reset is asserted, so u_cpu structurally never reboots). But
cpu_domain_rst_n = cpu_core_rst_n & pmu_cpu_rst_n_cpu_sync ALSO drops
to 0 for the PMU's own retention-hold reset (pmu_cpu_rst_n_cpu_sync),
not just for a genuine POR/PLL-unlock event. Wiring s_rst_n_i to
cpu_domain_rst_n therefore hard-FLUSHED this bridge's gray-code
pointers and any FIFO'd beats on every PMU power-down — a HARD FLUSH
is a fabricated, asynchronous, mid-transaction abort (see the "Reset
here is a HARD FLUSH" note in async_axi_fifo.sv's header), the exact
opposite of the retention the CPU side of the same power-cycle
guarantees. On resume, u_cpu's cache/AXI-arbiter FSM (rv32i_cache_arbiter.sv)
was still correctly waiting for the response to its outstanding
request (its own state WAS retained, per the CPU-side contract), but
the bridge that was supposed to deliver that response had already
discarded it — permanent deadlock (bead 2k8: instrumented evidence
showed dc_stall stuck at 1 forever post-resume with arvalid=0,
rvalid=0, i.e. the D-cache's AR handshake had already completed and it
was blocked waiting in the R phase for a beat the flush erased).

Fix: s_rst_n_i = cpu_core_rst_n, the CPU's TRUE reset (POR + PLL lock
only — see cpu_core_rst_n's declaration comment above), so the
bridge's CPU-side face retains exactly like the CPU it serves and
asserts ONLY for the same events that would legitimately reboot u_cpu
(never for a PMU power-cycle, which leaves cpu_core_rst_n at 1
throughout). u_cpu itself, u_apb_dbg_cdc's m_rst_n_i, and the ISO_CPU
clamp all correctly stay on cpu_domain_rst_n / pmu_cpu_iso_en_sync —
this fix touches ONLY u_cpu_axi_cdc's s_rst_n_i.
\end{textblock}

\subsection*{Before: logic        pll\_m\_psel, pll\_m\_penable, pll\_m\_pwrite;}

\begin{textblock}
Instantiated UNCONDITIONALLY (not generate-guarded on PLL_IMPL=="RNM"):
a guarded bridge would mean the default STUB regression exercises a
different netlist from the one being fixed, and u_apb_pll2_cdc set the
precedent of an unconditional bridge. In STUB mode core_clk == clk_i,
so this is a latency change only (one 4-phase RZ handshake added to
every APB_PLL access), not a functional one.
\end{textblock}

\subsection*{Before: logic        pll\_m\_psel, pll\_m\_penable, pll\_m\_pwrite;}

\begin{textblock}
Bootstrap safety (no self-brick risk): the bridge's source face is
held in reset by core_rst_n = rst_n_i & pll_locked, i.e. the bridge
cannot pass a transaction until the PLL has already locked once. This
cannot deadlock the PLL itself: pll_apb_regs' CONTROL register resets
to 32'h0000_0001 (pll_enable=1) and bit[0] is excluded from WMASK
(non-clearable — the GH #89 anti-brick fix, commit 110d7ac), so the
PLL always enables and locks with zero APB access required. Putting a
reset-gated bridge in front of the (already non-clearable) enable
register therefore cannot recreate that self-brick.
\end{textblock}

\subsection*{Before: logic cpu\_gated\_clk\_en;}

\begin{textblock}
ICG glitch-safety (rv32i_clock_gate.sv): en is latched transparent-low
and ANDed with clk, so `en` may only be sampled while clk is low —
any change to `en` while clk is high is invisible to gclk until the
next low phase, which is exactly what makes this ICG glitch-free
regardless of what drives `en`. cpu_gated_clk_en is a combinational AND
of two flop outputs both registered on cpu_core_clk (the same clk
driving this gate) — ~pmu_cpu_clk_dis_sync (u_cpu_clk_dis_sync) and
pmu_cpu_rst_n_cpu_sync (u_cpu_pmu_rst_sync) — so it is a clean,
synchronously-timed signal with no combinational glitch hazard of its
own, same as the single-signal case it replaces — safe to drive `en`
directly.
\end{textblock}

\subsection*{Before: logic pmu\_cpu\_rst\_n\_cpu\_sync;}

\begin{textblock}
"per-domain PMU reset" reuse cdc_reset_sync.sv's header anticipates.
\end{textblock}

\subsection*{Before: logic ext\_irq\_cpu\_sync, timer\_irq\_cpu\_sync;}

\begin{textblock}
level synchroniser cannot miss the assertion — see soc_top.sv header
"IRQ routing" note.
\end{textblock}

\end{relocated}

\section{\texttt{soc/sram\_controller.sv}}\label{tm:mod:sram-controller}

This section holds comment text that was in the source file. The text is unchanged. Each line has no leading comment marker and no indentation.

\begin{relocated}{rtl/soc/sram_controller.sv}

\subsection*{Before: \$fatal(1, "sram\_controller: SRAM\_SKY130 requires MEM\_WOR...}

\begin{textblock}
Elaboration system task (IEEE 1800-2017 §20.11): $fatal used
directly inside a generate scope -- NOT wrapped in an `initial`
block -- is evaluated at ELABORATION time, not simulation time.
The previous `initial $error(...)` form only ever fired if the
design was actually simulated, so `verilator --lint-only`
elaboration (which never runs simulation) silently let a wrong
MEM_WORDS through and truncate onto the macro's 10-bit
addr0/addr1. This form fails the build immediately, under lint
and synthesis elaboration alike, before a single clock is run.
\end{textblock}

\subsection*{Before: localparam int unsigned GROUP\_WORDS = 32;}

\begin{textblock}
  Stage 1 (input register slice): registers the accepted W beat's
  commit/data/strobe/index. This is the ONLY register anywhere in
  this pipeline whose D-input is a same-cycle function of an
  s_w*/crossbar signal — nothing past this point is.

  Stage 2 (per-group local registers): MEM_WORDS is split into
  NGROUPS contiguous groups of GROUP_WORDS words each. Each group has
  its OWN we/data registers, loaded ONLY when that group is the
  target of the stage-1 beat. grp_wdata_q[k] explicitly HOLDS
  (self-feedback) when group k is not selected, rather than always
  capturing wdata1_q — this makes every group's registers
  structurally distinct (a different D-input mux/compare term per
  group), which is what stops Yosys `opt_merge` from recognising the
  NGROUPS replicas as identical and collapsing them back into one
  high-fanout register. grp_we_q[k], by contrast, is a genuine
  1-cycle write-enable PULSE (explicitly driven to '0 when not
  selected, not held) — it must clear every cycle it isn't the
  target, or stage 3 would keep re-committing a stale write forever.

  Stage 3 (mem[] write): each group's registered we/data gates only
  that group's own GROUP_WORDS words. The fan-out of any one stage-2
  register is local to one group (<= GROUP_WORDS*SW endpoints), and
  the fan-out of any stage-1 register is local to NGROUPS compare
  terms (<= MEM_WORDS/GROUP_WORDS, e.g. 128 for the default
  MEM_WORDS=4096) — never the whole MEM_WORDS*DW array in one hop.
\end{textblock}

\subsection*{Before: typedef enum logic [0:0] \{R\_IDLE, R\_BUSY\} rstate\_e;}

\begin{textblock}
No RTL bypass/forwarding path is added into the write pipeline for
the in-flight window itself: a same-address forwarding compare would
have to reach into the per-group stage-2 registers from the read
FSM's held index, reintroducing a comparison that spans the very
group boundaries this fix was designed to keep separate. If a future
master needs a hardware-guaranteed same-cycle (or in-flight-window)
RAW bypass to this SRAM, prefer comparing the read's held index
against `idx1_q`/`grp1_sel`+`local1_onehot` (the pipeline's own
already-narrow per-stage state) over reintroducing any MEM_WORDS-wide
structure.
\end{textblock}

\subsection*{Before: logic r\_issue\_now;}

\begin{textblock}
see the Timing note above.
\end{textblock}

\subsection*{Before: logic             mb\_csb0\_m;}

\begin{textblock}
── SRAM macro instantiation ─────────────────────────────────────────────
Port 0 (RW, write-only here) <- write FSM; Port 1 (R-only) <- read FSM.
USE_POWER_PINS is intentionally not defined for this build (matches the
existing sky130_sram_1kbyte_1rw1r_32x256_8 cache-macro convention in
rtl/mem/rv32i_dcache.sv / rv32i_icache.sv): vccd1/vssd1 connectivity is
supplied by the LEF + PDN_MACRO_CONNECTIONS at PD time, not by RTL-level
power ports, so no ports exist to tie off here.
DFT Stage 4: mem_collar between the controller and the macro (functional side of the arm).
bist_en = 0 is a pass-through of every signal below, so the functional behaviour is unchanged.
\end{textblock}

\subsection*{Before: D-input structurally distinct from every other group's.}

\begin{textblock}
(holds its previous value) — see the header comment
above: this self-hold is what keeps this register's
\end{textblock}

\end{relocated}
% --- generated by PR7a: end








\chapter{CDC primitives}\label{tm:ch:mod-cdc-primitives}

This chapter describes how each module in this area works. The area holds the synchronizers and the asynchronous FIFOs.

% --- generated by PR7a: begin
\section{\texttt{soc/cdc/cdc\_gray\_fifo.sv}}\label{tm:mod:cdc-gray-fifo}

This section holds comment text that was in the source file. The text is unchanged. Each line has no leading comment marker and no indentation.

\begin{relocated}{rtl/soc/cdc/cdc_gray_fifo.sv}

\subsection*{Header}

\begin{textblock}
DEPTH=2 full flag (bead claude_verilog_test-3snb)
------------------------------------------------
Full means "the write pointer is exactly DEPTH ahead of the (synchronised)
read pointer", i.e. the two PTR_W-bit BINARY pointers are equal in the low
ADDR_W bits and differ in the wrap bit. In gray code that is: the top two
bits of the gray pointers are both inverted and every lower bit is equal —
the general (DEPTH >= 4) compare
    wr_gray_d == {~rd_sync[PTR_W-1:PTR_W-2], rd_sync[PTR_W-3:0]}.
At DEPTH=2, PTR_W = 2 and there are no lower bits: rd_sync[PTR_W-3:0] is
rd_sync[-1:0], not a legal range, so the compare degenerates to inverting the
WHOLE pointer:  wr_gray_d == ~rd_sync.  Checked exhaustively over the four
2-bit pointer values (binary b -> gray g = b ^ (b >> 1)):
    rd bin 00 (g 00) -> full at wr bin 10 (g 11) = ~00   ok
    rd bin 01 (g 01) -> full at wr bin 11 (g 10) = ~01   ok
    rd bin 10 (g 11) -> full at wr bin 00 (g 00) = ~11   ok
    rd bin 11 (g 10) -> full at wr bin 01 (g 01) = ~10   ok
and no other (wr, rd) pair satisfies it: the four gray codes are all
distinct, so wr_gray_d == ~rd_sync selects exactly one wr value per rd value.
Inverting only the MSB instead (the naive narrowing of the general form) is
WRONG at depth 2: it matches wr bin 11 against rd bin 00, i.e. it reports
full only after THREE words are in a two-entry memory (overflow).
Everything else — empty compare, pointer crossings, reset gating, storage —
is shared with DEPTH >= 4 unchanged. The depth-2 branch only selects the
compare operand (full_cmp_w below); no flop, port or instance name changes.
\end{textblock}

\subsection*{Before: if ((DEPTH < 2) || ((DEPTH \& (DEPTH - 1)) != 0)) begin :...}

\begin{textblock}
Elaboration-time guard: DEPTH must be a power of two and >= 2 (see
header). Generate-scope $fatal (not wrapped in `initial`) so it fires
under `verilator --lint-only` / synthesis elaboration and not only
under simulation — see rtl/soc/sram_controller.sv:110-126 for the
precedent and the reason `initial $fatal` alone is insufficient here.
\end{textblock}

\subsection*{Before: assign rd\_valid\_o = !empty\_q;}

\begin{textblock}
wr_ready_o is additionally gated by wr_rst_n_i itself: full_q's reset
value is 1'b0 ("not full"), so on its own !full_q would read 1 for the
entire window this domain is held in reset -- a write master could
then complete a handshake (wr_valid_i && wr_ready_o) whose beat is
silently dropped, since wr_bin_q is also held at its reset value and
wr_en's clocked branch never executes while !wr_rst_n_i. Gating on
wr_rst_n_i directly backpressures the attached master until this
domain's own reset has released, matching rd_valid_o's behaviour below
(empty_q already resets to 1'b1, so no equivalent gate is needed there
-- confirmed, not assumed, by inspection of the reset block above).
\end{textblock}

\subsection*{Before: the raw combinational condition — that is the intended C...}

\begin{textblock}
Registered comparisons (see header for why this is one cycle "behind"
\end{textblock}

\subsection*{Before: generate-selected, so the flop below is identical for ev...}

\begin{textblock}
(see "DEPTH=2 full flag" in the header). Only the compare operand is
\end{textblock}

\subsection*{Before: logic [WIDTH-1:0] mem\_q [DEPTH];}

\begin{textblock}
mem_q asynchronous-data-path warning in the header.
\end{textblock}

\end{relocated}

\section{\texttt{soc/cdc/cdc\_reset\_sync.sv}}\label{tm:mod:cdc-reset-sync}

This section holds comment text that was in the source file. The text is unchanged. Each line has no leading comment marker and no indentation.

\begin{relocated}{rtl/soc/cdc/cdc_reset_sync.sv}

\subsection*{Before: if (STAGES < 2) begin : g\_stages\_check}

\begin{textblock}
Elaboration-time guard: STAGES < 2 leaves no settling margin for the
deasserting edge. Generate-scope $fatal (not wrapped in `initial`) so
it fires under `verilator --lint-only` elaboration, not just
simulation — see rtl/soc/sram_controller.sv:110-126.
\end{textblock}

\subsection*{Before: (* magic\_cdc *)}

\begin{textblock}
Deliberately NOT written as a single always_ff looping/shifting over
an unpacked array or a self-referencing packed vector: this module's
rst_n_o is, by construction, always consumed as an ASYNC RESET by
flops in the destination domain (that is the entire purpose of a
reset synchronizer). Verilator's SYNCASYNCNET check has a false-
positive interaction with that combination specifically: a *single*
always_ff driving multiple array/vector elements via a for-loop, whose
result is later used as a downstream async reset, gets its array
elements misclassified as "used both synchronously and
asynchronously" (confirmed by isolated repro; cdc_sync uses the
identical array+for-loop shape and lints clean ONLY because its
outputs are consumed as data, never as a reset edge). One always_ff
per stage sidesteps the false positive entirely and is, if anything,
the more idiomatic description of a discrete synchroniser flop chain.
\end{textblock}

\subsection*{Before: logic rst\_n\_async;}

\begin{textblock}
scan mode (see the header). Scanmode low => rst_n_async === rst_n_i.
\end{textblock}

\subsection*{Before: rst\_n\_o === sync\_q[STAGES-1], the original output.}

\begin{textblock}
independent of the chain flops (see the header). Scanmode low =>
\end{textblock}

\end{relocated}

\section{\texttt{soc/cdc/cdc\_sync.sv}}\label{tm:mod:cdc-sync}

This section holds comment text that was in the source file. The text is unchanged. Each line has no leading comment marker and no indentation.

\begin{relocated}{rtl/soc/cdc/cdc_sync.sv}

\subsection*{Header}

\begin{textblock}
Why this converges where the reverted attempt did not: the reverted
version redrew its hold decision (`hold_sel`) ONLY inside `always @(d_i)`,
with nothing that ever cleared it back to 0 once d_i stopped moving — so a
hold could latch and stage0_d = sync_q[0] became a permanent self-loop
(measured: forcing it on gave async_axi_fifo TESTS=26 PASS=3 FAIL=23).
This version instead tracks `d_prev_q`, a plain register that is updated
EVERY destination-clock edge unconditionally (`d_prev_q <= d_i`, no
dependence on the hold decision). The hold decision is only even DRAWN
while `d_i !== d_prev_q`, and because d_prev_q always catches up to d_i on
the very next edge regardless of what was decided, that window can never
span more than one destination-clock cycle — there is no state here that
can persist once d_i has settled. This is the same guarantee OpenTitan's
own design provides via a differently-shaped mechanism (an unconditional
synchronous clear of `data_sel` on every posedge, racing a combinational,
event-gated re-arm on `src_data_i`) — both designs take BOTH the new
source value and the previously-captured value and evaluate the decision
in the DESTINATION clock's process, which is exactly the property the
reverted attempt lacked.
\end{textblock}

\subsection*{Before: if (STAGES < 2) begin : g\_stages\_check}

\begin{textblock}
Elaboration-time guard: STAGES < 2 defeats the purpose of a multi-flop
synchroniser (single-flop capture gives no metastability settling
time). Fails immediately under lint/synth elaboration, not just sim
(generate-scope $fatal, not wrapped in `initial` — see
rtl/soc/sram_controller.sv:110-126 for why `initial $fatal` alone is
not sufficient under `verilator --lint-only`).
\end{textblock}

\subsection*{Before: packed-vector parameter sized by WIDTH; no struct member...}

\begin{textblock}
Value of every stage while rst_n_i is asserted (see header). Plain
\end{textblock}

\subsection*{Before: instantiation site, so today's bit-exact deterministic r...}

\begin{textblock}
Simulation-only randomised settling-delay injection on stage 0 (see
header). Defeatable per instance; OFF by default at every existing
\end{textblock}

\subsection*{Before: — d\_prev\_q always catches up to d\_i on the very next edg...}

\begin{textblock}
self-clearing convergence (see "Why this converges" in the header)
\end{textblock}

\end{relocated}
% --- generated by PR7a: end








\chapter{DFT blocks}\label{tm:ch:mod-dft-blocks}

This chapter describes how each module in this area works. The area holds the JTAG TAP, the scan wrappers and the self-test blocks.

% --- generated by PR7a: begin
\section{\texttt{soc/dft/dft\_clk\_switch.sv}}\label{tm:mod:dft-clk-switch}

This section holds comment text that was in the source file. The text is unchanged. Each line has no leading comment marker and no indentation.

\begin{relocated}{rtl/soc/dft/dft_clk_switch.sv}

\subsection*{Header}

\begin{textblock}
Why the A term carries a third input, tck_i. clk_a_i is the GATED tck_i (clock
gate = latch + AND) and so falls AFTER tck_i: one delta cycle in zero-delay
simulation, the gate's delay in silicon. The A latch opens on !tck_i, so with
en_a_d already 1, en_a_q can rise while clk_a_i is still high, and
(clk_a_i & en_a_q) alone would emit a runt (zero width in RTL; Xcelium reports
it, and an event-ordering-dependent result is not acceptable in either
simulator). Guarding the A term with tck_i kills the race at the OUTPUT instead:
tck_i has already fallen in the delta in which the latch opens, so the A term
is 0 when en_a_q rises, whatever the gate's delay. Functionally the term is
redundant (clk_a_i high implies tck_i high, apart from that skew interval); it
costs one extra AND input on the clock path and is the output-side form of the
clock-gating hold check of en_a_q against the fall of clk_a_i.

Why the latch itself must stay tck-only. The first fix for this race opened
the A latch on `!tck_i && !clk_a_i`. That is correct in simulation but breaks
the DFT TAP clock rule: tools/dft/check_scan_clk_rst.py (make dft-scan-check)
requires every TAP-domain sequential cell (u_dft_ctrl.u_clk_sw) to be clocked
by the tck_i port ONLY, and a latch gated by clk_a_i traces through the TAP
clock gate u_tck_cg. The output guard keeps every sequential cell here on
tck_i or scan_clk_i. B needs no such term: clk_b_i is both the gated source and
the latch phase (one net, no skew between the two uses), so en_b_q can never
change while clk_b_i is high.

This differs from the usual flop-only enable-before-switch mux ON PURPOSE. A
flop retires the old selection only on an edge of the deselected clock, so a
stopped deselected clock deadlocks the switch (the reason section 15.5 first
rejected the idea: scan_clk_i is normally STOPPED by the tester). A latch is
open for the whole low phase, so a source that idles low can be switched OFF
with no clock edge at all. The price is that the latch needs the clock to be
low, which is the contract below.
\end{textblock}

\subsection*{Before: always\_latch begin}

\begin{textblock}
of the gated clk_a_i is closed by the tck_i term in the output, see the header.
\end{textblock}

\end{relocated}

\section{\texttt{soc/dft/mbist\_ctrl.sv}}\label{tm:mod:mbist-ctrl}

This section holds comment text that was in the source file. The text is unchanged. Each line has no leading comment marker and no indentation.

\begin{relocated}{rtl/soc/dft/mbist_ctrl.sv}

\subsection*{Before: 0: compare here from dout0\_i / dout1\_i; cmp\_fail*\_i are ...}

\begin{textblock}
1: collar-local compare (see header); dout0_i / dout1_i are unused, cmp_fail*_i and x_*_o are used.
\end{textblock}

\end{relocated}
% --- generated by PR7a: end








\chapter{PLL subsystem}\label{tm:ch:mod-pll-subsystem}

This chapter describes how each module in this area works. The area holds the PLL wrapper and its control registers.

% --- generated by PR7a: begin
\section{\texttt{soc/pll/pll\_subsystem.sv}}\label{tm:mod:pll-subsystem}

This section holds comment text that was in the source file. The text is unchanged. Each line has no leading comment marker and no indentation.

\begin{relocated}{rtl/soc/pll/pll_subsystem.sv}

\subsection*{Before: logic ref\_clk\_w;        // clk\_i, or scan\_clk\_i in scan ...}

\begin{textblock}
── DFT: test-clock and scan-reset bypass (see the header) ───────────────
\end{textblock}

\subsection*{Before: .clk\_i       (ref\_clk\_w),}

\begin{textblock}
Reference clock domain — see clock-domain note in module header.
\end{textblock}

\end{relocated}
% --- generated by PR7a: end





\chapter{Peripherals}\label{tm:ch:mod-peripherals}

This chapter describes how each module in this area works. The area holds the UART, the SPI, the timer and the other bus peripherals.

% --- generated by PR7b: begin
\section{\texttt{periph/dma\_engine.sv}}\label{tm:mod:dma-engine}

This section holds comment text that was in the source file. The text is unchanged. Each line has no leading comment marker and no indentation.

\begin{relocated}{rtl/periph/dma_engine.sv}

\subsection*{Header}

\begin{textblock}
Flat per-channel AXI ports (no SV interfaces) per Phase 5 RTL convention.
Lint target: verilator -Wall 0 errors 0 warnings.
\end{textblock}

\subsection*{Before: always\_ff @(posedge clk) begin}

\begin{textblock}
linebuf has NO reset and NO soft-reset term, deliberately (bead f7vs.15, Gate B).
It used to be written from inside the big FSM always_ff, i.e. under `if (!rst_n) ... else`
and `if (srst_pulse) ... else`, so rst_n (core_rst_n in soc_top) sat in the write-enable
cone of all MAX_BURST_BEATS*32 = 8192 flops. Gate B (Sky130 ss) measured that distribution at
98 stages / 30.8 ns. Giving the array its own always_ff keeps the same flop count (it was never
in the reset list) and removes core_rst_n from every one of those enables.

PROOF that writing linebuf during reset / soft-reset is unobservable.
linebuf is read in exactly one place: m_wdata = linebuf[wr_idx_q] in state S_W.
  1. S_W is entered only from S_AW (S_AW -> S_W on m_awready); S_AW is entered only from S_R
     on (m_rvalid && m_rresp == OKAY && m_rlast && rd_idx_q == beats_q-1), i.e. the one
     OKAY, exactly-beats_q-long burst. Every other path out of S_R is S_ERR or S_RDRAIN; S_RDRAIN
     exits only to S_ERR, and S_ERR goes to S_IDLE (halted_q), never to S_AW/S_W. So an S_W
     burst is always preceded, in the same transfer, by one complete, in-count run of S_R.
  2. Every run of S_R is preceded by S_CALC, which sets rd_idx_q <= 0 (the only other writer of
     rd_idx_q is the S_R increment and reset); the first S_R beat therefore writes linebuf[0],
     the next linebuf[1], and so on, one word per accepted m_rvalid, before S_AW can start.
     S_W then reads wr_idx_q = 0 .. beats_q-1 against those same indices. Because S_AW needs
     RLAST on beat beats_q-1 exactly (point 1), all beats_q words were written this transfer.
  3. Reset (rst_n low) and soft-reset (srst_pulse) both force state_q to S_IDLE. From S_IDLE
     the only way back to S_W is S_IDLE -> S_CALC -> S_AR -> S_R -> S_AW -> S_W, i.e. through
     point 2 again, which rewrites linebuf[0..beats-1] before any of it is read.
     Whatever the array held when reset/soft-reset hit -- including a word written on the very
     edge that state_q was forced to S_IDLE (state_q still reads S_R on that edge, so the new
     enable below fires where the old one was suppressed) -- is dead data.
  4. Error path: S_R writes the beat that carries the error (it is still a beat inside the
     burst, see below), then goes to S_ERR, or to S_RDRAIN when that beat was not RLAST. S_RDRAIN
     is NOT S_R, so the enable below is low for every drained beat: the drain writes nothing.
     The array content after an error is therefore dead data; the error path ends S_ERR ->
     S_IDLE with halted_q set, no S_W follows, and the next transfer restarts at S_CALC.
  5. In-range: the enable is state_q == S_R && m_rvalid and nothing else, with no beats_q
     comparator, because rd_idx_q < beats_q holds on every S_R cycle. S_CALC loads rd_idx_q = 0
     and beats_q >= 1; rd_idx_q increments only on an S_R beat that stays in S_R, and a beat
     stays in S_R only if it was error-free, which requires m_rlast == (rd_idx_q == beats_q-1),
     so a non-RLAST beat has rd_idx_q < beats_q-1 and the incremented index is <= beats_q-1. Any
     beat at index beats_q-1 leaves S_R (to S_AW, S_ERR or S_RDRAIN). An over-long burst's
     surplus beats are therefore taken in S_RDRAIN and cannot wrap rd_idx_q onto linebuf[0..].
     (rd_idx_q <= beats_q-1 <= MAX_BURST_BEATS-1 = 255 also keeps the 8-bit index in range.)
The enable below is therefore the minimum that preserves the old behaviour on every
reachable S_W read: (state_q == S_R) && m_rvalid. It has no rst_n and no srst_pulse term,
and the wdmo drain/in-range qualification is carried entirely by state_q and the
FSM's own transitions, so it adds no new term to the enable cone either.
\end{textblock}

\subsection*{Before: hw\_wen\_i  [REG\_CTRL]   <= 1'b0;}

\begin{textblock}
CTRL hardware-writeback flops (bead 0oc4): they MUST have a reset value. The register
bank samples hw_wen_i/hw_wdata_i on the first edge after rst_n deasserts; without a
reset these two held their power-up contents through the whole reset window, so in
silicon (or an X-propagating simulator) CTRL could load garbage there, including
CTRL.start (spurious DMA) and CTRL.soft_reset. Verilator zero-initialises flops, which
is why no ordinary simulation shows it (test_dma_ctrl_hw_flops_reset_* deposit garbage).
\end{textblock}

\end{relocated}

\section{\texttt{periph/gpio\_controller.sv}}\label{tm:mod:gpio-controller}

This section holds comment text that was in the source file. The text is unchanged. Each line has no leading comment marker and no indentation.

\begin{relocated}{rtl/periph/gpio_controller.sv}

\subsection*{Header}

\begin{textblock}
Input synchronisation: gpio_in_i is an asynchronous external pin bus. Per
rtl/soc/cdc/cdc_sync.sv's own scope warning, that primitive is valid
ONLY for WIDTH=1 or gray-coded buses -- NOT an arbitrary multi-bit binary
bus, because independent bits can resolve metastability on different
destination-clock edges and the sampled word could transiently equal a
value no single source-side cycle ever drove. GPIO pins carry no cross-bit
word coherency (each pin is an independent, unrelated external signal), so
this module instantiates N_PINS SEPARATE single-bit cdc_sync instances
-- one per pin -- rather than one WIDTH(N_PINS) instance, making that
absence of coherency explicit and self-documenting rather than accidental.
\end{textblock}

\subsection*{Header}

\begin{textblock}
This edge-sticky / level-live split follows ARM PL061 (GPIORIS is the raw
level in level mode; GPIOICR clears edge detection only). The alternative --
making level events sticky too -- was rejected: with the reset defaults
(IRQ_TYPE=0 level, IRQ_POL=0 active-low) and undriven pins reading low,
every pin's level condition is TRUE at reset, so a sticky raw register
would latch to all-ones on the first cycle out of reset and stay there
until software cleared a register it had not yet configured.
\end{textblock}

\subsection*{Header}

\begin{textblock}
Lint target: verilator -Wall -Wno-IMPORTSTAR 0 errors 0 warnings.
\end{textblock}

\subsection*{Before: if (N\_PINS == 0 || N\_PINS > 32) begin : g\_n\_pins\_check}

\begin{textblock}
Elaboration-time guard (IEEE 1800-2017 §20.11): $fatal directly inside a
generate scope, not wrapped in `initial`, fires at ELABORATION time --
under `verilator --lint-only` too, which never runs simulation and would
otherwise silently let an illegal N_PINS through. Mirrors
rtl/soc/cdc/cdc_sync.sv's g_stages_check and the rationale recorded
at rtl/soc/sram_controller.sv:110-126.
N_PINS == 0 is rejected as well as > 32: a zero-pin instance would make
every [N_PINS-1:0] part-select below a [-1:0] range, and PIN_MASK's
>> (32 - N_PINS) a >> 32 (all zeroes), so the module would be both
illegal and useless.
\end{textblock}

\subsection*{Before: localparam logic [31:0] PIN\_MASK = 32'((\{32\{1'b1\}\} >> (3...}

\begin{textblock}
Implemented-pin mask: only the low N_PINS bits are SW-writable.
Idiom mirrored verbatim from interrupt_controller.sv:77 (GH #87-safe
width-cast: N_PINS==32 must not truncate to 0).
\end{textblock}

\subsection*{Before: logic [WORDW-1:0] addr\_word\_w;}

\begin{textblock}
=========================================================================
GPIO_IRQ_CLR write-snoop — one-cycle clear mask.
Address-compare zero-extends addr_word to 32 b before comparing (GH #87
truncation-safe idiom, mirrored verbatim from apb4_register_bank.sv) so
that N_REGS == 2**WORDW can never wrap a valid index to a false match.
=========================================================================
\end{textblock}

\end{relocated}

\section{\texttt{periph/i2c\_bit\_engine.sv}}\label{tm:mod:i2c-bit-engine}

This section holds comment text that was in the source file. The text is unchanged. Each line has no leading comment marker and no indentation.

\begin{relocated}{rtl/periph/i2c_bit_engine.sv}

\subsection*{Header}

\begin{textblock}
  * The release wait is not extra time on top of four ticks: the 3 clk it takes the released SCL
    to come back through the 2-FF synchroniser (2 flops + the state register) are COUNTED INSIDE
    the high phase, by shortening HI1 by that latency (an uncompensated fourth tick would give a
    4N + 3 period -- 3 clk slow). HI1 >= 1 clk needs CLKDIV_eff >= SYNC_STAGES + 1, which is
    exactly the CLKDIV_MIN elaboration guard in i2c_controller.sv.
  * The K skew moves floor(N/8) clk from the high phase into the SCL-low phase WITHOUT changing the
    period. A symmetric 2N / 2N split cannot meet the Fast-mode tLOW (1.3 us) at 400 kHz: two ticks
    of 0.625 us give 1.25 us. Each phase length is checked EXHAUSTIVELY for every CLKDIV_eff
    3..65535 at the exact maximum-rate clock of each mode (f_clk = 4*N*f_max; a slower bus at the
    same f_clk only lengthens both phases). Standard mode (tLOW 4.7 / tHIGH 4.0 us at <= 100 kHz)
    is met for every N. Fast mode (1.3 / 0.6 us at <= 400 kHz) is met for every N EXCEPT
    N = 4..7 (K = 0) and N = 13..15 (K = 1): there floor(N/8) is too small and tLOW lands at
    1.25..1.30 us -- i.e. only when f_clk is 6.4..11.2 MHz or 20.8..24 MHz AND the bus is run at
    exactly 400 kHz. None of this SoC's clocks (40 MHz Sky130, 571 MHz ASAP7, 100 MHz test) hits
    it; at those N run the bus a little slower. For N >= 16 the low time is >= 0.52 * T and the
    high time >= 0.46 * T. (ceil(N/8) would shrink the failing set to N = 4 alone; it changes the
    CLKDIV 4 low time and so needs test_i2c_clkdiv_clamp_to_min's "+2" assertion updated.)
\end{textblock}

\subsection*{Header}

\begin{textblock}
Lint target: verilator -Wall -Wno-IMPORTSTAR 0 errors 0 warnings.
\end{textblock}

\subsection*{Before: localparam int unsigned SYNC\_LAT   = SYNC\_STAGES + 1;   ...}

\begin{textblock}
Tick counter: reloaded on every state change and on every tick, so each interval state lasts
reload+1 clk from entry. In the wait states it free-runs at one tick per CLKDIV_eff+1 clk,
clocking the timeout (the wait states always reload the plain CLKDIV_eff, so I2C_TIMEOUT stays
in whole ticks). The reload value is chosen by the state being ENTERED:
  S_BIT_L2  : CLKDIV_eff + K              (SCL-low phase carries the duty-cycle skew)
  S_BIT_HI1 : CLKDIV_eff - (SYNC_STAGES+1) - K   (the sync wait + skew come out of the high phase)
  S_ST_LO / S_RS_HI / S_STP_HI : 2*N - 1 = 2*CLKDIV_eff + 1   (TICKS_COND = 2 ticks)
  S_STP_FREE: 3*N - 1 = 3*CLKDIV_eff + 2                      (TICKS_FREE = 3 ticks)
  otherwise : CLKDIV_eff
K = floor(N / 8), N = CLKDIV_eff + 1. The first two are >= 0 for CLKDIV_eff >= SYNC_STAGES + 1
(the CLKDIV_MIN guard); 18 b arithmetic cannot overflow (max 3*0xFFFF + 2 = 0x2FFFF).
\end{textblock}

\end{relocated}

\section{\texttt{periph/i2c\_controller.sv}}\label{tm:mod:i2c-controller}

This section holds comment text that was in the source file. The text is unchanged. Each line has no leading comment marker and no indentation.

\begin{relocated}{rtl/periph/i2c_controller.sv}

\subsection*{Header}

\begin{textblock}
Measured (100 MHz clk): CLKDIV 249 -> 1000 clk = 100.0 kHz (tLOW 5.31 us, tHIGH 4.69 us); CLKDIV 62
-> 252 clk = 396.8 kHz (tLOW 1.33 us, tHIGH 1.19 us).
\end{textblock}

\subsection*{Header}

\begin{textblock}
Prescaler minimum (explicit elaboration guard): the engine samples the SYNCHRONISED lines, which
lag the pins by SYNC_STAGES clk, and the oe register adds one more. An SCL-low interval is two
ticks and the SDA-to-sample distance is at least two ticks, so a tick must comfortably exceed
that latency. CLKDIV_MIN (default 3 -> tick = 4 clk, 2x margin over the 4-clk worst case) must be
>= SYNC_STAGES + 1; a smaller override FAILS ELABORATION (g_clkdiv_min_check), and CLKDIV below
CLKDIV_MIN is clamped in hardware. With the four-phase bit (see "Bit timing") the SAME bound is
also what keeps the shortened S_BIT_HI1 reload (CLKDIV_eff - (SYNC_STAGES+1) - K) non-negative:
at CLKDIV_eff = SYNC_STAGES + 1 it is exactly 0 (a 1-clk phase, K = 0), so the minimum is unchanged
by the fourth phase and is tight, not conservative. The SCL-low interval is now 2N + K >= two ticks
and the SDA-to-sample distance 2N + K - 3 >= two ticks, so the original latency argument holds.
The bit engine receives the clamped value and does not re-check it.
The same lesson as spi_controller.sv's SPI_CLK_DIV >= 7, but enforced rather than documented (SPI
leaves 0..1 unsafe).
\end{textblock}

\subsection*{Header}

\begin{textblock}
Lint target: verilator -Wall -Wno-IMPORTSTAR 0 errors 0 warnings.
\end{textblock}

\subsection*{Before: logic [WORDW-1:0] addr\_word\_w;}

\begin{textblock}
-------------------------------------------------------------------------
APB address decode (zero-extended to 32 b before comparing -- GH #87-safe idiom).
-------------------------------------------------------------------------
\end{textblock}

\subsection*{Before: the clamp constant inside the 16-bit CLKDIV field.}

\begin{textblock}
The prescaler minimum must exceed the synchroniser latency (see header). The upper bound keeps
\end{textblock}

\subsection*{Before: logic [15:0] div\_eff\_w;}

\begin{textblock}
Clamp: a tick must exceed the synchroniser latency (see header).
\end{textblock}

\end{relocated}

\section{\texttt{periph/interrupt\_controller.sv}}\label{tm:mod:interrupt-controller}

This section holds comment text that was in the source file. The text is unchanged. Each line has no leading comment marker and no indentation.

\begin{relocated}{rtl/periph/interrupt_controller.sv}

\subsection*{Header}

\begin{textblock}
Lint target: verilator -Wall -Wno-IMPORTSTAR 0 errors 0 warnings.
\end{textblock}

\subsection*{Before: logic [N\_SOURCES-1:0] irq\_src\_sync\_q;   // synchronised ...}

\begin{textblock}
=========================================================================
2-FF synchroniser for irq_src_i (CDC hardening)
irq_src_i originates from peripheral IRQ lines. In soc_top every source
is currently synchronous to this clock (core_clk, or the gated copy of it
that clocks the GPU), so for today's SoC the chain is a two-cycle
retiming stage; it is kept as a real synchroniser so a source that moves
to another clock domain stays safe.

N_SOURCES independent single-bit cdc_sync instances (STAGES=2, reset
0), not one WIDTH(N_SOURCES) instance: the sources are independent
level-held lines with no cross-bit coherency requirement, and
cdc_sync.sv's scope warning forbids a multi-bit binary bus. Same
pattern as gpio_controller.sv (g_in_sync). Behaviour vs the previous
hand-written chain: same 2-cycle latency and same reset value (0); the
reset is now asserted asynchronously by the primitive (see the
cdc_sync header) where it was synchronous here.
=========================================================================
\end{textblock}

\end{relocated}

\section{\texttt{periph/pwm\_controller.sv}}\label{tm:mod:pwm-controller}

This section holds comment text that was in the source file. The text is unchanged. Each line has no leading comment marker and no indentation.

\begin{relocated}{rtl/periph/pwm_controller.sv}

\subsection*{Header}

\begin{textblock}
Lint target: verilator -Wall -Wno-IMPORTSTAR 0 errors 0 warnings.
\end{textblock}

\subsection*{Before: if (N\_CH == 0 || N\_CH > 8) begin : g\_n\_ch\_check}

\begin{textblock}
Elaboration-time guard (IEEE 1800-2017 §20.11): $fatal directly inside a generate scope,
not wrapped in `initial`, fires at ELABORATION time -- under `verilator --lint-only` too.
Mirrors rtl/soc/cdc/cdc_sync.sv's g_stages_check and gpio_controller.sv:137-139.
\end{textblock}

\subsection*{Before: localparam logic [31:0] CH\_MASK = 32'((\{32\{1'b1\}\} >> (32...}

\begin{textblock}
Implemented-channel mask, MAP_CH-wide. Idiom mirrored from interrupt_controller.sv:78 /
gpio_controller.sv:160 (GH #87-safe width-cast: MAP_CH==32 case never arises here, but the
idiom is kept identical for consistency).
\end{textblock}

\subsection*{Before: logic [WORDW-1:0] addr\_word\_w;}

\begin{textblock}
=========================================================================
PWM_IRQ_CLR write-snoop — one-cycle clear mask.
Address-compare zero-extends addr_word to 32 b before comparing (GH #87
truncation-safe idiom, mirrored verbatim from apb4_register_bank.sv /
gpio_controller.sv:277-289) so N_REGS == 2**WORDW can never wrap a valid
index to a false match.
=========================================================================
\end{textblock}

\subsection*{Before: verilator lint\_off UNUSEDSIGNAL */}

\begin{textblock}
The expanded strobe MUST land in a named net before being sliced.

Writing `strb_expand(pstrb)[N_CH-1:0]` -- a part-select applied directly to a function
call -- is legal SystemVerilog and Verilator accepts it, but it is NOT legal
Verilog-2005. sv2v passes the construct through verbatim rather than lowering it, and
yosys then rejects the generated file with "syntax error, unexpected '['". That breaks
the CDC flow AND both PD flows (Sky130 and ASAP7 are both sv2v-fronted), while no cocotb
suite can catch it, since Verilator consumes the SystemVerilog and never sees sv2v output.

The 32-bit temp costs a local UNUSEDSIGNAL waiver for bits [31:N_CH] -- deliberately
preferred over a portability break. strb_expand stays byte-for-byte identical to
apb4_register_bank.sv / gpio_controller.sv rather than being narrowed to N_CH.
\end{textblock}

\end{relocated}

\section{\texttt{periph/spi\_controller.sv}}\label{tm:mod:spi-controller}

This section holds comment text that was in the source file. The text is unchanged. Each line has no leading comment marker and no indentation.

\begin{relocated}{rtl/periph/spi_controller.sv}

\subsection*{Header}

\begin{textblock}
Lint target: verilator -Wall -Wno-IMPORTSTAR 0 errors 0 warnings.
\end{textblock}

\subsection*{Before: typedef enum logic [0:0] \{}

\begin{textblock}
=========================================================================
SPI shift engine FSM

Mode table (CPOL, CPHA):
  CPOL=0, CPHA=0: idle SCLK=0; sample rising, shift falling  (Mode 0)
  CPOL=0, CPHA=1: idle SCLK=0; shift rising,  sample falling (Mode 1)
  CPOL=1, CPHA=0: idle SCLK=1; sample falling, shift rising  (Mode 2)
  CPOL=1, CPHA=1: idle SCLK=1; shift falling,  sample rising (Mode 3)

In all modes: the "leading edge" is the first toggle of SCLK.
  CPHA=0: sample on leading edge (active-to-idle direction for CPOL=1,
          idle-to-active for CPOL=0), shift on trailing edge.
  CPHA=1: shift on leading edge, sample on trailing edge.

We track edge count (0..15 for 8 full cycles = 16 edges).
sclk_q toggles on each sclk_edge; we derive sample/shift events from
whether it's an odd or even edge count, combined with CPHA.
=========================================================================
\end{textblock}

\subsection*{Before: logic spi\_miso\_sync\_q;}

\begin{textblock}
spi_miso_i CDC synchroniser (bead claude_verilog_test-a2g, GH #95
cdc_snitch finding spi-miso-unsynchronised-sample). spi_miso_i is a
board-level asynchronous pin; sampling it directly into rx_shift_q /
rx_push_data was benign only by construction (SPI master mode drives
SCLK itself), not clean by inspection. Route it through the standard
2-FF synchroniser instead of hand-rolling flops.

Timing impact: cdc_sync adds SYNC_STAGES (default 2) fabric-clock cycles of latency
between a spi_miso_i transition and its visibility to the shift
engine. This is transparent as long as the external slave's MISO has
already been stable for >= SYNC_STAGES cycles by the sample edge -- true for any
half-period (SPI_CLK_DIV+1) comfortably larger than 2, which is what
this module's own external-slave guidance already requires (see the
SPI_CLK_DIV >= 7 recommendation above). Loopback mode is unaffected:
miso_loopback never routes through spi_miso_i/spi_miso_sync_q.
KNOWN GAP: SPI_CLK_DIV == 0 (half-period = 1 clock) and SPI_CLK_DIV == 1
(half-period = 2 clocks) are NOT safe for a real external slave with
this synchroniser in place -- the 2-cycle sync latency meets or exceeds
the half-period, so the sampled bit can be stale by a full bit
position. Flagged, not silently retimed; see bead
claude_verilog_test-a2g close-out notes. Raising SYNC_STAGES raises the
minimum safe half-period by the same amount: the SPI_CLK_DIV >= 7
guidance assumes SYNC_STAGES = 2, so keep the half-period
(SPI_CLK_DIV + 1 clocks) comfortably above SYNC_STAGES.
\end{textblock}

\end{relocated}

\section{\texttt{periph/trng.sv}}\label{tm:mod:trng}

This section holds comment text that was in the source file. The text is unchanged. Each line has no leading comment marker and no indentation.

\begin{relocated}{rtl/periph/trng.sv}

\subsection*{Header}

\begin{textblock}
Lint target: verilator -Wall -Wno-IMPORTSTAR 0 errors 0 warnings.
\end{textblock}

\subsection*{Before: logic [WORDW-1:0] addr\_word\_w;}

\begin{textblock}
-------------------------------------------------------------------------
APB address decode (zero-extended to 32 b before comparing -- GH #87-safe idiom).
-------------------------------------------------------------------------
\end{textblock}

\end{relocated}

\section{\texttt{periph/trng\_lfsr\_entropy.sv}}\label{tm:mod:trng-lfsr-entropy}

This section holds comment text that was in the source file. The text is unchanged. Each line has no leading comment marker and no indentation.

\begin{relocated}{rtl/periph/trng_lfsr_entropy.sv}

\subsection*{Header}

\begin{textblock}
Lint target: verilator -Wall -Wno-IMPORTSTAR 0 errors 0 warnings.
\end{textblock}

\end{relocated}

\section{\texttt{periph/uart\_controller.sv}}\label{tm:mod:uart-controller}

This section holds comment text that was in the source file. The text is unchanged. Each line has no leading comment marker and no indentation.

\begin{relocated}{rtl/periph/uart_controller.sv}

\subsection*{Header}

\begin{textblock}
Lint target: verilator -Wall -Wno-IMPORTSTAR 0 errors 0 warnings.
\end{textblock}

\subsection*{Before: logic [WORDW-1:0] apb\_word\_addr;}

\begin{textblock}
=========================================================================
APB access-phase snoop — detect UART_TX write and UART_RX read commits.

APB4 access phase completes when psel & penable & pready are all 1.
pready is driven 1'b1 combinatorially by apb4_register_bank (zero wait
states), so the access phase is always exactly one cycle wide.  No capture
registers are needed — the word address, pwrite, and pstrb are all stable
during the setup and access phases by APB4 protocol.

Word address from the byte address:
  addr_word = paddr[ADDR_W-1:2]

TX push: APB write ACCESS to UART_TX with >=1 strobe asserted.
  A zero-strobe write (pstrb==4'h0) is a legal APB no-op and must not push.

RX pop:  APB read ACCESS completing at UART_RX word address.
  APB is single-transfer; the access pulse is one cycle → no double-pop risk.
=========================================================================
\end{textblock}

\subsection*{Before: logic [15:0] baud\_cnt\_q;   // oversample period down-cou...}

\begin{textblock}
=========================================================================
A4 — 16x oversample tick generator (shared by TX and RX engines)

UART_BAUD register (D = regs_o[REG_UART_BAUD][15:0]) sets the oversample
period: 1 os_tick = (D+1) system clocks.
1 serial bit = 16 os_ticks = 16*(D+1) system clocks.
When D=0 → os_tick every clock; bit period = 16 clocks.

os_phase counts 0..15; wraps every 16 os_ticks → one bit_tick per wrap.
TX FSM uses bit_tick; RX FSM uses os_tick + os_phase for mid-bit sampling.
TX and RX share this generator so loopback stays exactly bit-aligned.
=========================================================================
\end{textblock}

\subsection*{Before: always\_ff @(posedge clk) begin}

\begin{textblock}
The shared phase counter free-runs: it times the TX engine only.  The RX
engine keeps its own phase counter (rx_phase_q, below) that it re-aligns on
every START edge, so receiving never disturbs the transmitted bit timing
and TX stays exactly 16 os_ticks per bit (bead shm5).
\end{textblock}

\subsection*{Before: typedef enum logic [2:0] \{}

\begin{textblock}
=========================================================================
A4 — RX engine (16x oversample, serial receiver; re-syncs every frame)

Receives from uart_rx_i (through the cdc_sync synchroniser) or from the
internal tx_serial when loopback is enabled (same clock domain — no
synchroniser needed).

Own phase counter (bead shm5):
  rx_phase_q counts os_ticks 0..15 and is reset to 0 on the os_tick at
  which RX_IDLE sees the line low (the START falling edge, seen up to one
  os_tick late).  It is independent of the TX phase counter os_phase_q, so
  re-aligning the receiver never stretches or shortens a transmitted bit.
  Every frame therefore starts from a fresh alignment and a baud error
  cannot accumulate from one byte to the next.

Mid-bit sampling with majority vote (A3-integrated):
  Each bit is voted 2-of-3 on rx_phase_q = 6, 7, 8.  Because the phase
  counter restarts AT the detection os_tick, rx_phase_q == k is seen k+1
  os_ticks after it; phases 6/7/8 are therefore 7/8/9 os_ticks after the
  detected edge, centred on 8 ticks, the bit mid-point (the true edge is
  0..1 os_tick before the detected one).
  START bit: a majority-high vote (not genuinely low mid-bit) is a noise
             glitch; it is rejected one os_tick after the vote completes.
  Data bits: voted values are latched at the end of the bit (phase 15).
  STOP bit:  voted on phases 6, 7 and the live line at phase 8; the frame
             is closed AT that mid-point: framing error if the vote is 0,
             byte pushed, RX_IDLE re-armed.  The next START may be
             accepted from there on, so frames sent back-to-back with a
             single stop bit and a bit-clock error of a few percent keep
             their framing.  The os_ticks between that mid-point and the
             nominal end of the stop bit absorb the error: the receiver
             is back in IDLE before even a fast sender's next START.

Loopback alignment:
  TX advances on bit_tick; the START it drives is seen by RX one os_tick
  later and the receiver aligns to it.  The shared os_tick generator keeps
  the two engines on the same tick grid.
=========================================================================
\end{textblock}

\subsection*{Before: logic uart\_rx\_sync\_q;}

\begin{textblock}
2-FF synchroniser on the asynchronous uart_rx_i pad (CDC hardening).
Idle line is high (8N1), so the synchroniser must reset to 1: a reset
value of 0 would present a false START bit to the receiver FSM.
cdc_sync with RESET_VAL = 1 resets both stages to 1 (no inverters).
Loopback path uses the internal tx_serial (same clock domain) and
bypasses this synchroniser. Same 2-cycle latency as the previous
hand-written chain; the reset is now asserted asynchronously by the
primitive (see the cdc_sync header) where it was synchronous here.
SYNC_STAGES > 2 is safe for the receiver: the RX FSM aligns its
oversample phase to the *synchronised* start edge and every later data
edge is delayed by the same constant, so the sample points keep their
position relative to the bits; only the data-valid latency grows.
\end{textblock}

\end{relocated}

\section{\texttt{periph/watchdog\_timer.sv}}\label{tm:mod:watchdog-timer}

This section holds comment text that was in the source file. The text is unchanged. Each line has no leading comment marker and no indentation.

\begin{relocated}{rtl/periph/watchdog_timer.sv}

\subsection*{Header}

\begin{textblock}
Lint target: verilator -Wall -Wno-IMPORTSTAR 0 errors 0 warnings.
\end{textblock}

\subsection*{Before: logic [WORDW-1:0] addr\_word\_w;}

\begin{textblock}
-------------------------------------------------------------------------
APB address decode (zero-extended to 32 b before comparing -- GH #87-safe idiom, as in
apb4_register_bank.sv / gpio_controller.sv / pwm_controller.sv).
-------------------------------------------------------------------------
\end{textblock}

\end{relocated}
% --- generated by PR7b: end








\chapter{Crypto accelerator}\label{tm:ch:mod-crypto-accelerator}

This chapter describes how each module in this area works. The area holds the AES and SHA cores and their register interface.

% --- generated by PR7b: begin
\section{\texttt{periph/aes128\_core.sv}}\label{tm:mod:aes128-core}

This section holds comment text that was in the source file. The text is unchanged. Each line has no leading comment marker and no indentation.

\begin{relocated}{rtl/periph/aes128_core.sv}

\subsection*{Header}

\begin{textblock}
NO RUNTIME-INDEXED MUX (bead ma7 acceptance criterion; the proven Synlig OPT_MUXTREE miscompile).
  * The S-box is `function sbox` = `unique case (b)` over 256 CONSTANT labels returning
    literals -- NOT an indexed read of a 256-entry localparam array.
  * The round constant is `rcon_q <= xtime(rcon_q)`, reset 8'h01, advanced once per round: no
    table and no select.
  * ShiftRows is fixed wiring; MixColumns and SubBytes use constant slices only; the one
    select left in the datapath is the 1-bit mix_en 2:1 (round 10 skips MixColumns).

KEY SCHEDULE: ON-THE-FLY. rk_q is ONE 128-bit register holding the current round key, advanced
once per round: w0' = w0 ^ SubWord(RotWord(w3)) ^ {rcon,24'h0}, w1' = w0'^w1, w2' = w1'^w2,
w3' = w2'^w3. Precomputed-and-stored (11 round keys = 1408 flops) is FORBIDDEN here twice over:
the flop count, and because fetching round key r is exactly the forbidden runtime-indexed read.

MIXCOLUMNS per column a0..a3, t = a0^a1^a2^a3: b_i = a_i ^ t ^ xtime(a_i ^ a_(i+1 mod 4)).
Identity: a0 ^ t ^ xtime(a0^a1) = 2a0 ^ 3a1 ^ a2 ^ a3. 4 xtimes per column (16 total) against 8 for
the naive 2a_i ^ 3a_(i+1) ^ a_(i+2) ^ a_(i+3) expansion.
\end{textblock}

\subsection*{Header}

\begin{textblock}
Fmax note: SBOX_PARALLEL=16 critical path is S-box -> ShiftRows wiring -> 2 XOR levels of
MixColumns -> AddRoundKey in one cycle; comfortable at 40 MHz Sky130. AREA (20 x 256x8) is the
Gate A question, not the period.
\end{textblock}

\subsection*{Header}

\begin{textblock}
Lint target: verilator -Wall -Wno-IMPORTSTAR 0 errors 0 warnings.
\end{textblock}

\subsection*{Before: function automatic logic [7:0] sbox(input logic [7:0] b)...}

\begin{textblock}
=========================================================================
Pure functions (declaration before use: yosys-slang strict mode, bead q7n)
=========================================================================
\end{textblock}

\subsection*{Before: function automatic logic [7:0] sbox(input logic [7:0] b)...}

\begin{textblock}
AES S-box: 256 CONSTANT case labels, never an indexed array read (bead ma7).
\end{textblock}

\subsection*{Before: s\_q     <= block\_i \^{} key\_i;    // initial AddRoundKey (w...}

\begin{textblock}
PRE-LOAD WHILE IDLE (bead f7vs.15, Gate B). The whitening value and the round-1
key/counters are loaded on EVERY idle cycle; start_i only advances state_q. The
values captured on the start edge are exactly the ones the old `if (start_i)` load
captured (same block_i ^ key_i, key_i, 8'h01, 4'd1 sampled at that same edge), so the
cycle count and every result are unchanged: 11 / 41 clk from the start edge. What
changes is the cone: start_i is the APB decode (paddr -> start_pulse_w) and used to
gate the D-input of all 128 + 128 + 8 + 4 datapath flops, and yosys' sbox ROM
inference (proc_rom + memory_dff) merged the S-box address flops, so that decode
also sat in front of an 8-bit S-box network. Now it fans out to state_q only.
None of s_q/rk_q/rcon_q/round_q is observable while idle (busy_o, done_set_o need
state_q == A_ROUND; dout_o is consumed only with done_set_o), and they hold no secret
that block_i/key_i do not already hold in the parent's msg_q/key_q.
\end{textblock}

\end{relocated}

\section{\texttt{periph/crypto\_accel.sv}}\label{tm:mod:crypto-accel}

This section holds comment text that was in the source file. The text is unchanged. Each line has no leading comment marker and no indentation.

\begin{relocated}{rtl/periph/crypto_accel.sv}

\subsection*{Header}

\begin{textblock}
NO NEW RUNTIME-INDEXED MUX (bead ma7 acceptance criterion). The S-box and the SHA K table are
constant-label case functions, Rcon is an xtime() chain, the message schedule a shift register;
see the sub-core headers. In this file every KEY / DIN / DOUT / DIGEST word is addressed by a
constant index inside a constant-bound unrolled loop.
\end{textblock}

\subsection*{Header}

\begin{textblock}
REGISTER-BANK PATTERN. hw_wen = 1 every cycle <=> the word is a LIVE MIRROR of internal state
(STATUS, IRQ_STAT). hw_wen = a pulse <=> the bank word IS the storage element and must hold
between events (DOUT, DIGEST, IV3). The second form is why this design has no result registers
at all. STATUS and IRQ_STAT are written every cycle with the value the state flops (state_q,
done_q, key_rej_q, key_seen_q) are ABOUT TO TAKE, so after every clock edge each field of those
words equals the corresponding state flop's new value, with no lag: a transfer that completes at
edge N is already reflected in a read whose ACCESS phase follows it. That is the whole guarantee
-- STATUS == f(state flops) after every edge -- and it implies nothing about any other register.
\end{textblock}

\subsection*{Header}

\begin{textblock}
Lint target: verilator -Wall -Wno-IMPORTSTAR 0 errors 0 warnings.
\end{textblock}

\subsection*{Before: logic [WORDW-1:0] addr\_word\_w;}

\begin{textblock}
=========================================================================
Snoop front end (one copy, shared by every snoop). Only one APB transfer completes per
cycle, so no two snoops can ever fire together.
Address compares zero-extend addr_word_w to 32 b before comparing (GH #87 idiom, mirrored
from apb4_register_bank.sv) so N_REGS == 2**WORDW can never wrap a valid index to a match.
=========================================================================
\end{textblock}

\subsection*{Before: logic scan\_s1\_q, scan\_s2\_q, scan\_prev\_q;}

\begin{textblock}
-------------------------------------------------------------------------
Scan-mode edge detect -> wipe_w (bead j41m.6). See the SCAN WIPE note in the header.
2-FF synchroniser + one history flop, deliberately UNRESET (they track scan_mode_i through
reset so that a reset can neither create nor hide an edge).
-------------------------------------------------------------------------
\end{textblock}

\subsection*{Before: assign scan\_wipe\_done\_o = \~{}((scan\_mode\_i \^{} scan\_prev\_q) ...}

\begin{textblock}
Wipe-done (bead j41m.17): every stage of the edge-detect pipeline, and the raw level, agree
with the history flop. It falls combinationally when scan_mode_i changes (before any clk
edge), stays 0 while the change walks s1 -> s2 -> prev, and rises just after the edge that
loads scan_prev_q -- the SAME edge at which the wipe_w-qualified clears take effect. A change
of scan_mode_i back before the wipe completes keeps it low until the pipeline has flushed.
In functional mode (scan_mode_i = 0 for >= 3 clk edges) it is a constant 1.
\end{textblock}

\subsection*{Before: if (wipe\_w) begin}

\begin{textblock}
SCAN WIPE (bead j41m.6): the bank words that hold key- or message-derived results are
cleared through the hardware-write port; this block comes last so it overrides every
completion / increment write above in the same cycle. IV is not secret but is counter
state that scan must not leave live either.
\end{textblock}

\end{relocated}

\section{\texttt{periph/sha256\_core.sv}}\label{tm:mod:sha256-core}

This section holds comment text that was in the source file. The text is unchanged. Each line has no leading comment marker and no indentation.

\begin{relocated}{rtl/periph/sha256_core.sv}

\subsection*{Header}

\begin{textblock}
NO RUNTIME-INDEXED MUX (bead ma7 acceptance criterion; the proven Synlig OPT_MUXTREE miscompile).
  * The K constants are `function sha_k` = `unique case (t)` over 64 CONSTANT labels returning
    literals -- NOT `K_TABLE[t]`. A case over constants collapses to a 6-input constant
    function per output bit; an indexed array read is the $pmux form that miscompiled in ma7.
  * The message schedule is a 16 x 32 SHIFT REGISTER with constant taps only (below).
  * Sigma/sigma are constant rotations and shifts, i.e. wiring.

W ROLLING SCHEDULE. Window convention: at the start of round t, w_q[j] = W[t+j]; the round
consumes w_q[0]. Every round cycle
    w_q[j]  <= w_q[j+1]                                       j = 0..14
    w_q[15] <= sigma1(w_q[14]) + w_q[9] + sigma0(w_q[1]) + w_q[0]
which is exactly W[t+16] = s1(W[t+14]) + W[t+9] + s0(W[t+1]) + W[t]. W[16] is therefore produced
during round 0. Rounds 48..63 compute words that are never consumed; they are left to run
(gating them saves nothing and adds a condition).
\end{textblock}

\subsection*{Header}

\begin{textblock}
  The 65-cycle variant -- folding the eight final adds onto the round-63 next-state -- is
  REJECTED: the round path already carries Sigma0, Sigma1, Ch, Maj and two 32-bit adder chains,
  and stacking another 32-bit add there is the one place in this peripheral that could move
  Fmax. Do not "optimise" it in later.
\end{textblock}

\subsection*{Header}

\begin{textblock}
Lint target: verilator -Wall -Wno-IMPORTSTAR 0 errors 0 warnings.
\end{textblock}

\subsection*{Before: FIPS 180-4 Sec.4.2.2 round constants: 64 CONSTANT case l...}

\begin{textblock}
=========================================================================
Pure functions (declaration before use: yosys-slang strict mode, bead q7n)
=========================================================================
\end{textblock}

\end{relocated}
% --- generated by PR7b: end








\chapter{NPU}\label{tm:ch:mod-npu}

This chapter describes how each module in this area works. The area holds the INT8 array, the weight memory and the register interface.

% --- generated by PR7b: begin
\section{\texttt{npu/npu\_mac\_array.sv}}\label{tm:mod:npu-mac-array}

This section holds comment text that was in the source file. The text is unchanged. Each line has no leading comment marker and no indentation.

\begin{relocated}{rtl/npu/npu_mac_array.sv}

\subsection*{Header}

\begin{textblock}
Lint target: verilator --lint-only -Wall -Wno-IMPORTSTAR 0 errors 0 warnings.
\end{textblock}

\end{relocated}

\section{\texttt{npu/npu\_top.sv}}\label{tm:mod:npu-top}

This section holds comment text that was in the source file. The text is unchanged. Each line has no leading comment marker and no indentation.

\begin{relocated}{rtl/npu/npu_top.sv}

\subsection*{Header}

\begin{textblock}
NO NEW RUNTIME-INDEXED MUX. The grid is genvar-structural, the weight address goes to a macro,
the requantiser's lane select is a one-hot AND-OR, and the engine state is one-hot shift
registers. The only read multiplexer is apb4_register_bank's, already in every netlist.

BANK PATTERN. hw_wen = 1 every cycle <=> a LIVE MIRROR written with the NEXT-STATE value
(STATUS, AOUT, IRQ_STAT), so after every edge each field equals the state it mirrors, with no
lag. WADDR is the one pulse-written word (the bank word IS the address counter). The bank cannot
protect a register dynamically, so "a WADDR write while busy is rejected" is done by gating
penable into the bank for exactly that transfer; the plan's WMASK for WADDR is unchanged.
\end{textblock}

\subsection*{Header}

\begin{textblock}
Lint target: verilator --lint-only -Wall -Wno-IMPORTSTAR 0 errors 0 warnings.
\end{textblock}

\end{relocated}

\section{\texttt{npu/npu\_weight\_mem.sv}}\label{tm:mod:npu-weight-mem}

This section holds comment text that was in the source file. The text is unchanged. Each line has no leading comment marker and no indentation.

\begin{relocated}{rtl/npu/npu_weight_mem.sv}

\subsection*{Header}

\begin{textblock}
THERE IS NO BEHAVIOURAL ARRAY IN THIS FILE, so no define can ever infer flops. The
sram_controller.sv two-way shape (macro / flat `logic [31:0] mem [..]`) is deliberately NOT
copied: its flat arm is the 32,768-flop structure bead rvb proved non-terminating in the ASAP7
post-CTS resizer, and the NPU would have been a second one.
\end{textblock}

\subsection*{Header}

\begin{textblock}
NO RUNTIME-INDEXED MUX (bead ma7). The weight address goes to a macro, not to a mux. The only
selector is the 4:1 bank read mux of the two banked arms, written as a tree of 2:1 ternaries on
the REGISTERED bank select with constant bank indices -- never bank_dout[sel_q]. Bank decode for
csb0 is a constant compare per generate instance.
\end{textblock}

\subsection*{Header}

\begin{textblock}
Lint target: verilator --lint-only -Wall -Wno-IMPORTSTAR 0 errors 0 warnings (default,
+define+SRAM_SKY130 and +define+SRAM_ASAP7 builds).
\end{textblock}

\end{relocated}
% --- generated by PR7b: end








\chapter{FPGA wrappers}\label{tm:ch:mod-fpga-wrappers}

This chapter describes how each module in this area works. The area holds the memory models and the wrappers for the FPGA build.

% --- generated by PR7c: begin
\section{\texttt{fpga/fpga\_sdp\_ram\_be.sv}}\label{tm:mod:fpga-sdp-ram-be}

This section holds comment text that was in the source file. The text is unchanged. Each line has no leading comment marker and no indentation.

\begin{relocated}{rtl/fpga/fpga_sdp_ram_be.sv}

\subsection*{Header}

\begin{textblock}
Inference rules: synchronous read, no reset on array or output register, byte-write written as a for-loop
over constant byte slices of one array element, array indexed only by waddr_i / raddr_i. No `initial`.
\end{textblock}

\end{relocated}

\section{\texttt{fpga/fpga\_sram\_1rw.sv}}\label{tm:mod:fpga-sram-1rw}

This section holds comment text that was in the source file. The text is unchanged. Each line has no leading comment marker and no indentation.

\begin{relocated}{rtl/fpga/fpga_sram_1rw.sv}

\subsection*{Header}

\begin{textblock}
Block-RAM inference rules honoured: synchronous read, NO reset on the array or the read register, ONE
address shared by read and write, the array indexed only by that address (no runtime-indexed mux of
flops). No `initial`, no vendor primitive: the same file is valid for yosys synth_xilinx / synth_ecp5 /
vendor tools. Read-during-write cannot occur (one port, read and write are exclusive by web0).

The "hold on a write" semantic is a read-enable (csb0 low AND web0 high) on the output register, which is
the no-change / enable behaviour block RAMs have natively; if a toolchain cannot map it, it falls back to
a read-first block RAM plus a feedback mux and the report must say so.
\end{textblock}

\end{relocated}
% --- generated by PR7c: end
